% 第10章 定积分——题目解答
% 仅包含原章练习、思考题、参考题；普通已解例题不重复。

\section*{第10章 定积分——题目解答}

\subsection*{10.1.3 练习题}

\paragraph{练习题 1}
\begin{xsolution}
不可积。事实上，在任意 $x_0\in(0,1)$ 处，由有理数与无理数的稠密性，均可取有理数列 $r_n\to x_0$ 与无理数列 $s_n\to x_0$。于是
\[
  f(r_n)\to x_0(1-x_0)>0,
  \qquad
  f(s_n)=0,
\]
故 $f$ 在每个 $x_0\in(0,1)$ 处都不连续。另一方面，在 $0,1$ 两点有
\[
  \abs{f(x)}\leqslant x(1-x)\to0,
\]
所以 $f$ 只在端点连续。

因此 $f$ 的不连续点集为 $(0,1)$，不是零测度集。由 Lebesgue 判别定理（命题 10.1.6），$f\notin R[0,1]$。
\end{xsolution}

\paragraph{练习题 2}
\begin{xsolution}
第一个问题的答案是否定的。取 $[a,b]=[0,1]$，令 $f$ 为 Riemann 函数（Thomae 函数）
\[
  f(x)=
  \begin{cases}
    1/q,&x=p/q\in[0,1]\text{，且 }(p,q)=1,\\
    0,&x\text{ 为无理数},
  \end{cases}
\]
则 $f\in R[0,1]$，且 $f([0,1])\subseteq[0,1]$。再令
\[
  g(y)=
  \begin{cases}
    1,&y=0,\\
    0,&y\ne0.
  \end{cases}
\]
函数 $g$ 只在 $0$ 处与零函数不同，故 $g\in R[0,1]$。但
\[
  (g\circ f)(x)=
  \begin{cases}
    0,&x\text{ 为有理数},\\
    1,&x\text{ 为无理数},
  \end{cases}
\]
是 Dirichlet 型函数，在每一点都不连续，因而不可积。

第二个问题的答案也是否定的。仍在 $[0,1]$ 上令
\[
  f(x)=g(x)=
  \begin{cases}
    1,&x\text{ 为有理数},\\
    0,&x\text{ 为无理数}.
  \end{cases}
\]
则 $f,g$ 都不可积，而 $f(x)$ 只取 $0,1$ 两值，且 $g(0)=g(1)=1$，所以
\[
  (g\circ f)(x)\equiv1,
\]
反而是可积函数。
\end{xsolution}

\paragraph{练习题 3}
\begin{xsolution}
有如下关系：
\[
  f\in R[a,b]\Longrightarrow \abs{f}\in R[a,b]
  \Longleftrightarrow f^2\in R[a,b].
\]
第一步是因为绝对值函数连续，连续函数与 Riemann 可积函数的复合仍可积。若 $\abs f\in R[a,b]$，则
\[
  f^2=(\abs f)^2\in R[a,b].
\]
反之，若 $f^2\in R[a,b]$，则 $f^2\geqslant0$，而平方根函数在 $[0,+\infty)$ 上连续，所以
\[
  \abs f=\sqrt{f^2}\in R[a,b].
\]

反向推出 $f$ 一般不成立。例如在任意非退化区间上令
\[
  f(x)=
  \begin{cases}
    1,&x\text{ 为有理数},\\
    -1,&x\text{ 为无理数},
  \end{cases}
\]
则 $f$ 处处不连续，从而不可积，但
\[
  \abs f\equiv1,\qquad f^2\equiv1
\]
均可积。
\end{xsolution}

\paragraph{练习题 4}
\begin{xsolution}
先设 $f,g$ 只在有限个点 $x_1,\ldots,x_N$ 上不同。令 $h=g-f$。则 $h$ 只在有限点可能非零，因此 $h\in R[a,b]$ 且
\[
  \int_a^b h(x)\,\dd x=0.
\]
于是
\[
  g=f+h\in R[a,b],
  \qquad
  \int_a^b g=\int_a^b f.
\]

若只把“有限个点”放宽为“几乎处处相等”，则一般不能推出相同结论。反例取
\[
  f(x)\equiv0,
  \qquad
  g(x)=
  \begin{cases}
    1,&x\in\QQ,\\
    0,&x\notin\QQ,
  \end{cases}
  \qquad x\in[a,b].
\]
有理数集为零测度集，所以 $f=g$ 几乎处处；但 $f\in R[a,b]$，而 $g$ 是 Dirichlet 型函数，在每一点都不连续，故 $g\notin R[a,b]$。

如果另加条件 $g\in R[a,b]$，则积分值仍然相同。事实上 $h=g-f\in R[a,b]$ 且 $h=0$ 几乎处处，因此 $\int_a^b h=0$，从而 $\int_a^b f=\int_a^b g$。
\end{xsolution}

\paragraph{练习题 5}
\begin{xsolution}
所求积分为 $0$。

由于 $f\in R[a,b]$，由 Lebesgue 判别定理，$f$ 的连续点在 $[a,b]$ 中稠密。设 $x_0$ 是一个连续点。若 $f(x_0)>0$，由连续性存在包含 $x_0$ 的小区间 $(\alpha,\beta)$，使其中恒有 $f(x)>0$。于是任意 $x_1,x_2\in(\alpha,\beta)$ 都满足
\[
  f(x_1)f(x_2)>0,
\]
与题设矛盾。同理不可能有 $f(x_0)<0$。故 $f$ 在每个连续点上都等于 $0$。

因此 $f=0$ 几乎处处，而 $f$ Riemann 可积，故
\[
  \boxed{\int_a^b f(x)\,\dd x=0}.
\]
\end{xsolution}

\paragraph{练习题 6}
\begin{xsolution}
因为 $g\in R[a,b]$，故 $g$ 有界，且其不连续点集为零测度集。题设给出
\[
  m\leqslant g(x)\leqslant M.
\]
函数 $f\in C[m,M]$ 在紧区间上连续，从而一致连续。若 $x_0$ 是 $g$ 的连续点，则
\[
  g(x)\to g(x_0)\quad(x\to x_0)
\]
立即推出
\[
  f(g(x))\to f(g(x_0)),
\]
所以 $f\circ g$ 在 $x_0$ 连续。因此
\[
  D(f\circ g)\subseteq D(g),
\]
其中 $D(\cdot)$ 表示不连续点集。右边为零测度集，而且 $f\circ g$ 有界。由 Lebesgue 判别定理，
\[
  f\circ g\in R[a,b].
\]
\end{xsolution}

\paragraph{练习题 7}
\begin{xsolution}
由 $1/f$ 在 $[a,b]$ 上有界，存在常数 $C>0$，使
\[
  \abs*{\frac1{f(x)}}\leqslant C,
\]
故
\[
  \abs{f(x)}\geqslant c:=\frac1C>0.
\]
函数 $\varphi(y)=1/y$ 在集合
\[
  (-\infty,-c]\cup[c,+\infty)
\]
上连续，而且在 $f([a,b])$ 上一致连续。于是 $f$ 的每个连续点也是 $1/f$ 的连续点。因此
\[
  D(1/f)\subseteq D(f).
\]
因 $f\in R[a,b]$，$D(f)$ 为零测度集；同时 $1/f$ 已知有界。由 Lebesgue 判别定理即得
\[
  \frac1f\in R[a,b].
\]
\end{xsolution}

\paragraph{练习题 8}
\begin{xsolution}
设 $f$ 的所有间断点组成收敛数列 $\{x_n\}$；若极限点本身未列入数列，把它补入也只增加一个点。于是 $f$ 的不连续点集至多可列，因而为零测度集。题设又给出 $f$ 在 $[a,b]$ 上有界，所以由 Lebesgue 判别定理
\[
  f\in R[a,b].
\]
\end{xsolution}

\paragraph{练习题 9}
\begin{xsolution}
对每个 $x_0\in[a,b]$，题设给出
\[
  \lim_{x\to x_0}f(x)=0.
\]
先证 $f$ 有界。对每个 $x_0$，存在邻域 $U_{x_0}$，使在 $U_{x_0}\cap[a,b]$ 内（除去 $x_0$ 本身）有 $\abs{f(x)}<1$。加上点值 $f(x_0)$ 后，$f$ 在每个这样的邻域上有界。利用 $[a,b]$ 的紧致性取有限子覆盖，即知 $f$ 在整个 $[a,b]$ 上有界。

再看非零点集。对任意 $n\in\NN_+$，令
\[
  E_n=\{x\in[a,b]\mid \abs{f(x)}\geqslant1/n\}.
\]
$E_n$ 不可能有聚点：若 $x_k\in E_n$ 且 $x_k\to x_0$、$x_k\ne x_0$，则由题设应有 $f(x_k)\to0$，矛盾。紧区间中的无限集合必有聚点，所以每个 $E_n$ 都是有限集。于是
\[
  \{x\mid f(x)\ne0\}=\bigcup_{n=1}^\infty E_n
\]
至多可列。

若 $f(x_0)=0$，由极限也为 $0$ 可知 $f$ 在 $x_0$ 连续；若 $f(x_0)\ne0$，则该点是不连续点。因此 $f$ 的不连续点集至多可列。由 Lebesgue 判别定理，$f\in R[a,b]$。又 $f=0$ 几乎处处，故
\[
  \boxed{\int_a^b f(x)\,\dd x=0}.
\]
\end{xsolution}

\paragraph{练习题 10}
\begin{xsolution}
Riemann 函数在任意有界区间上有界；由本章正文已知，它恰在有理点处不连续，而在无理点处连续。有理数集至多可列，因而为零测度集。故由 Lebesgue 判别定理，Riemann 函数在每个有界区间 $[a,b]$ 上 Riemann 可积。

并且它在所有无理数点取值为 $0$，故也有
\[
  \int_a^b f(x)\,\dd x=0.
\]
\end{xsolution}

\paragraph{练习题 11}
\begin{xsolution}
可以。事实上，对连续函数，这种等距分割右端点和确实足以确定定积分。

设
\[
  S_n=\frac{b-a}{n}\sum_{i=1}^n
  f\left(a+\frac{i}{n}(b-a)\right).
\]
由于 $f\in C[a,b]$，$f$ 在 $[a,b]$ 上一致连续。令等距分割的步长为
\[
  \Delta_n=\frac{b-a}{n}.
\]
任取从属于该分割的介点 $\xi_i$，则
\[
  \abs*{S_n-\sum_{i=1}^n f(\xi_i)\Delta_n}
  \leqslant
  \sum_{i=1}^n
  \abs*{f\left(a+i\Delta_n\right)-f(\xi_i)}\Delta_n.
\]
当 $n$ 充分大时，每个 $\abs{a+i\Delta_n-\xi_i}\leqslant\Delta_n$ 都很小；由一致连续性，右端可任意小。另一方面，连续函数 Riemann 可积，任意介点 Riemann 和都趋于 $\int_a^b f$。因此
\[
  \lim_{n\to\infty}S_n=\int_a^b f(x)\,\dd x.
\]
若题设又给出 $S_n\to I$，由极限唯一性即得
\[
  I=\int_a^b f(x)\,\dd x.
\]
\end{xsolution}

\paragraph{练习题 12}
\begin{xsolution}
由于 $f,g\in R[a,b]$，二者都有界，且乘积 $fg\in R[a,b]$。取常数 $A,B>0$，使
\[
  \abs{f(x)}\leqslant A,
  \qquad
  \abs{g(x)}\leqslant B.
\]
记
\[
  S_P=\sum_{k=1}^n f(x_k)\sin(g(x_k)\Delta x_k),
  \qquad
  T_P=\sum_{k=1}^n f(x_k)g(x_k)\Delta x_k.
\]
利用
\[
  \abs{\sin u-u}\leqslant\frac{\abs u^3}{6},
\]
有
\begin{align*}
  \abs{S_P-T_P}
  &\leqslant\frac{A}{6}\sum_{k=1}^n
    \abs{g(x_k)}^3(\Delta x_k)^3\\
  &\leqslant\frac{AB^3}{6}\norm P^2\sum_{k=1}^n\Delta x_k\\
  &=\frac{AB^3(b-a)}6\norm P^2.
\end{align*}
故当 $\norm P\to0$ 时，$S_P-T_P\to0$。而 $T_P$ 是 $fg$ 的右端点 Riemann 和，所以
\[
  T_P\to\int_a^b f(x)g(x)\,\dd x.
\]
于是
\[
  S_P\to\int_a^b f(x)g(x)\,\dd x.
\]
按极限定义，对任意 $\varepsilon>0$ 可取 $\delta>0$，使 $\norm P<\delta$ 时题中不等式成立。
\end{xsolution}

\subsection*{10.2.1 思考题}

\paragraph{思考题（1）}
\begin{xsolution}
取区间 $[-1,1]$，令
\[
  f(x)=x,\qquad g(x)=x.
\]
这时 $g$ 在区间上变号，而且
\[
  \int_{-1}^1 g(x)\,\dd x=0,
  \qquad
  \int_{-1}^1 f(x)g(x)\,\dd x
  =\int_{-1}^1x^2\,\dd x=\frac23.
\]
若积分第一中值定理在没有保号条件时仍成立，则右端应为
\[
  f(\xi)\int_{-1}^1g(x)\,\dd x=0,
\]
与 $2/3$ 矛盾。因此 $g$ 的保号性不能任意删去。
\end{xsolution}

\paragraph{思考题（2）}
\begin{xsolution}
取区间 $[0,1]$，令
\[
  f(x)\equiv1,\qquad g(x)=x(1-x).
\]
函数 $g$ 不是单调函数，并且
\[
  g(0)=g(1)=0,
  \qquad
  \int_0^1f(x)g(x)\,\dd x=\int_0^1x(1-x)\,\dd x=\frac16.
\]
而积分第二中值定理的结论若成立，则对某个 $\xi\in[0,1]$ 应有
\[
  \int_0^1f(x)g(x)\,\dd x
  =g(0)\int_0^\xi f+g(1)\int_\xi^1 f=0,
\]
矛盾。因此单调性条件也不能任意删去。
\end{xsolution}

\subsection*{10.2.4 练习题}

\paragraph{练习题 1}
\begin{xsolution}
取
\[
  g(x)=(x-a)(b-x)f(x).
\]
则 $g\in C[a,b]$ 且 $g(a)=g(b)=0$。由题设
\[
  0=\int_a^b f(x)g(x)\,\dd x
  =\int_a^b(x-a)(b-x)f^2(x)\,\dd x.
\]
被积函数在 $[a,b]$ 上连续且非负，所以它处处为 $0$。在 $x\in(a,b)$ 上有 $(x-a)(b-x)>0$，故 $f(x)=0$。再由连续性得 $f(a)=f(b)=0$。因此
\[
  f\equiv0.
\]
\end{xsolution}

\paragraph{练习题 2}
\begin{xsolution}
反设 $P\not\equiv0$。由 $f\geqslant0$ 且 $\int_a^b f>0$，例题 10.2.1 保证存在子区间 $[c,d]\subseteq[a,b]$ 与 $\mu>0$，使
\[
  f(x)\geqslant\mu\qquad(x\in[c,d]).
\]
非零多项式只有有限个零点，故可在 $(c,d)$ 中取一个更小的闭区间 $[c_1,d_1]$，其中 $P$ 无零点。由连续性存在 $\nu>0$，使
\[
  P^2(x)\geqslant\nu\qquad(x\in[c_1,d_1]).
\]
于是
\[
  \int_a^bP^2(x)f(x)\,\dd x
  \geqslant\int_{c_1}^{d_1}P^2(x)f(x)\,\dd x
  \geqslant\mu\nu(d_1-c_1)>0,
\]
与题设矛盾。因此只能有 $P\equiv0$。
\end{xsolution}

\paragraph{练习题 3}
\begin{xsolution}
令
\[
  h(x)=f(x)+f(-x).
\]
则 $h$ 是 $[-1,1]$ 上的连续偶函数，因此可在题设中取 $g=h$，得到
\[
  \int_{-1}^1 f(x)h(x)\,\dd x=0.
\]
把积分作代换 $x\mapsto-x$，又有
\[
  \int_{-1}^1 f(-x)h(x)\,\dd x=0.
\]
两式相加：
\[
  \int_{-1}^1 h^2(x)\,\dd x=0.
\]
连续非负函数积分为零只能处处为零，故
\[
  h(x)=f(x)+f(-x)\equiv0.
\]
因此 $f(-x)=-f(x)$，即 $f$ 为奇函数。
\end{xsolution}

\paragraph{练习题 4}
\begin{xsolution}
极限为 $0$。任取 $0<\delta<1$，分拆积分：
\[
  0\leqslant\int_0^1(1-x^2)^n\,\dd x
  \leqslant
  \int_0^\delta1\,\dd x+
  \int_\delta^1(1-\delta^2)^n\,\dd x.
\]
所以
\[
  0\leqslant\int_0^1(1-x^2)^n\,\dd x
  \leqslant\delta+(1-\delta)(1-\delta^2)^n.
\]
先令 $n\to\infty$，得上极限不超过 $\delta$；再令 $\delta\to0+$，即得
\[
  \boxed{\lim_{n\to\infty}\int_0^1(1-x^2)^n\,\dd x=0}.
\]
\end{xsolution}

\paragraph{练习题 5}
\begin{xsolution}
给出一个自然的推广如下：设 $f\in C[a,b]$，满足
\[
  0\leqslant f(x)\leqslant1,
\]
且集合
\[
  E=\{x\in[a,b]\mid f(x)=1\}
\]
为零测度集，则
\[
  \lim_{n\to\infty}\int_a^b f^n(x)\,\dd x=0.
\]

证明如下。给定 $\varepsilon>0$，可用有限个开区间覆盖紧集 $E$，并使这些区间与 $[a,b]$ 的交的总长度小于 $\varepsilon/2$。把这些开区间的并记为 $U$。紧集
\[
  K=[a,b]\setminus U
\]
与 $E$ 不交，因此由连续性存在 $q<1$，使
\[
  f(x)\leqslant q\qquad(x\in K).
\]
于是
\[
  0\leqslant\int_a^b f^n(x)\,\dd x
  \leqslant\frac\varepsilon2+(b-a)q^n.
\]
当 $n$ 充分大时第二项小于 $\varepsilon/2$，故积分小于 $\varepsilon$。于是极限为 $0$。

例题 10.2.4 对应 $f(x)=\sin x$、$[a,b]=[0,\pi/2]$；这时 $f=1$ 只发生在端点 $\pi/2$，故满足上述条件。
\end{xsolution}

\paragraph{练习题 6}
\begin{xsolution}
\textbf{解 1}\quad 由题设
\[
  \sin^n x_n=\frac2\pi I_n,
  \qquad
  I_n:=\int_0^{\pi/2}\sin^n x\,\dd x.
\]
先给出 $I_n$ 的一个不呈指数衰减的下界。对充分大的 $n$，在
\[
  x\in\left[\frac\pi2-\frac1{\sqrt n},\frac\pi2\right]
\]
上写 $y=\frac\pi2-x$，则 $0\leqslant y\leqslant1/\sqrt n$，且
\[
  \sin x=\cos y\geqslant1-\frac{y^2}{2}
  \geqslant1-\frac1{2n}.
\]
因此
\[
  I_n\geqslant\frac1{\sqrt n}
  \left(1-\frac1{2n}\right)^n.
\]
右边括号的 $n$ 次方趋于 $\ee^{-1/2}>0$，故存在 $c>0$，使充分大的 $n$ 有
\[
  \sin^n x_n=\frac2\pi I_n\geqslant\frac{c}{\sqrt n}.
\]

若 $x_n$ 不趋于 $\pi/2$，则存在 $\delta>0$ 与子列 $n_k$，使
\[
  x_{n_k}\leqslant\frac\pi2-\delta.
\]
令 $q=\sin(\frac\pi2-\delta)<1$，则
\[
  \sin^{n_k}x_{n_k}\leqslant q^{n_k},
\]
而 $q^{n_k}$ 比 $n_k^{-1/2}$ 衰减快，和上面的下界矛盾。因此
\[
  \boxed{\lim_{n\to\infty}x_n=\frac\pi2}.
\]

\textbf{另解}\quad 记
\[
  I_n=\int_0^{\pi/2}\sin^n x\,\dd x.
\]
题设给出
\[
  \sin x_n=\left(\frac2\pi I_n\right)^{1/n}
  =\left(\frac2\pi\right)^{1/n}I_n^{1/n}.
\]
由例题 10.2.5 对 $f(x)=\sin x$ 的结论，
\[
  \lim_{n\to\infty}I_n^{1/n}
  =\max_{0\leqslant x\leqslant\pi/2}\sin x=1,
\]
而 $(2/\pi)^{1/n}\to1$，故 $\sin x_n\to1$。因为 $x_n\in[0,\pi/2]$ 且 $\sin x$ 在该区间严格增加，遂有
\[
  \boxed{\lim_{n\to\infty}x_n=\frac\pi2}.
\]
\end{xsolution}

\paragraph{练习题 7}
\begin{xsolution}
记
\[
  K_h(x)=\frac{h}{h^2+x^2},\qquad h>0.
\]
把积分写为
\[
  \int_{-1}^1K_h(x)f(x)\,\dd x
  =f(0)\int_{-1}^1K_h(x)\,\dd x
  +\int_{-1}^1K_h(x)(f(x)-f(0))\,\dd x.
\]
第一项中
\[
  \int_{-1}^1K_h(x)\,\dd x
  =2\arctan\frac1h\longrightarrow\pi.
\]

只需证明余项趋于 $0$。给定 $\varepsilon>0$，由 $f$ 在 $0$ 连续，取 $\delta\in(0,1)$，使 $\abs{x}<\delta$ 时
\[
  \abs{f(x)-f(0)}<\varepsilon.
\]
又令 $M=\max_{[-1,1]}\abs{f(x)-f(0)}$。则
\begin{align*}
  \abs*{\int_{-1}^1K_h(x)(f(x)-f(0))\,\dd x}
  &\leqslant
  \varepsilon\int_{-\delta}^{\delta}K_h(x)\,\dd x
  +2M\int_\delta^1\frac{h}{x^2}\,\dd x\\
  &\leqslant \pi\varepsilon+\frac{2Mh}{\delta}.
\end{align*}
先令 $h\to0+$，再令 $\varepsilon\to0+$，余项为 $0$。故
\[
  \boxed{\lim_{h\to0+}\int_{-1}^1\frac{h}{h^2+x^2}f(x)\,\dd x=\pi f(0)}.
\]
\end{xsolution}

\paragraph{练习题 8（1）}
\begin{xsolution}
设 $M=\max_{[0,1]}\abs{f(x)}$，则
\[
  \abs*{\int_0^1x^nf(x)\,\dd x}
  \leqslant M\int_0^1x^n\,\dd x
  =\frac{M}{n+1}\to0.
\]
故极限为
\[
  \boxed{0}.
\]
\end{xsolution}

\paragraph{练习题 8（2）}
\begin{xsolution}
写成
\[
  n\int_0^1x^nf(x)\,\dd x
  =f(1)\frac{n}{n+1}
  +n\int_0^1x^n(f(x)-f(1))\,\dd x.
\]
第一项趋于 $f(1)$。对第二项，给定 $\varepsilon>0$，由 $f$ 在 $1$ 连续，可取 $\delta\in(0,1)$，使 $1-\delta<x\leqslant1$ 时
\[
  \abs{f(x)-f(1)}<\varepsilon.
\]
若 $M=\max_{[0,1]}\abs{f(x)-f(1)}$，则
\begin{align*}
  n\int_0^1x^n\abs{f(x)-f(1)}\,\dd x
  &\leqslant nM\int_0^{1-\delta}x^n\,\dd x
  +n\varepsilon\int_{1-\delta}^1x^n\,\dd x\\
  &\leqslant M(1-\delta)^{n+1}+\varepsilon.
\end{align*}
令 $n\to\infty$ 后再令 $\varepsilon\to0$，第二项趋于 $0$。因此
\[
  \boxed{\lim_{n\to\infty}\int_0^1nx^nf(x)\,\dd x=f(1)}.
\]
\end{xsolution}

\paragraph{练习题 9}
\begin{xsolution}
直接计算
\[
  \int_0^{a_n}x^n\,\dd x=\frac{a_n^{n+1}}{n+1}\to2.
\]
所以
\[
  a_n^{n+1}=2(n+1)(1+\littleo(1)).
\]
取对数：
\[
  (n+1)\ln a_n=\ln(n+1)+\ln2+\littleo(1).
\]
除以 $n+1$，右端趋于 $0$，故 $\ln a_n\to0$，从而
\[
  \boxed{\lim_{n\to\infty}a_n=1}.
\]
\end{xsolution}

\paragraph{练习题 10}
\begin{xsolution}
记
\[
  I_n=\int_0^{\pi/2}\frac{\sin^2nt}{\sin t}\,\dd t.
\]
利用恒等式
\[
  \sin^2nt-\sin^2(n-1)t
  =\sin( (2n-1)t)\sin t,
\]
得到
\begin{align*}
  I_n-I_{n-1}
  &=\int_0^{\pi/2}\sin((2n-1)t)\,\dd t\\
  &=\frac1{2n-1}.
\end{align*}
并且 $I_1=1$，所以
\[
  I_n=\sum_{k=1}^n\frac1{2k-1}
  =H_{2n}-\frac12H_n,
\]
其中 $H_n=\sum_{k=1}^n1/k$。由 $H_n=\ln n+\bigO(1)$，
\[
  I_n=\frac12\ln n+\bigO(1).
\]
因此
\[
  \boxed{\lim_{n\to\infty}\frac{I_n}{\ln n}=\frac12}.
\]
\end{xsolution}


\subsection*{10.3.1 思考题}

\paragraph{思考题（1）}
\begin{xsolution}
可取
\[
  f(x)=\operatorname{sgn}x,\qquad x\in[-1,1],
\]
并任意补定 $f(0)$。该函数只有一个跳跃间断点，所以 Riemann 可积；但导函数具有 Darboux 性质，而 $\operatorname{sgn}x$ 在 $0$ 两侧分别取负值与正值，却不具有介值性，因此它不可能是任何函数的导函数，即没有原函数。
\end{xsolution}

\paragraph{思考题（2）}
\begin{xsolution}
令
\[
  F(x)=
  \begin{cases}
    x^2\sin\dfrac1{x^2},&x\ne0,\\
    0,&x=0.
  \end{cases}
\]
则 $F$ 在 $[-1,1]$ 上处处可导，且
\[
  F'(0)=0,
\]
而当 $x\ne0$ 时
\[
  F'(x)=2x\sin\frac1{x^2}-\frac2x\cos\frac1{x^2}.
\]
第二项在 $x\to0$ 时无界，因此导函数 $F'$ 在 $[-1,1]$ 上无界。Riemann 可积函数必有界，所以 $F'$ 有原函数 $F$，但 $F'$ 不可积。
\end{xsolution}

\subsection*{10.3.3 练习题}

\paragraph{练习题 1（1）}
\begin{xsolution}
令
\[
  G(x)=\int_0^x t\sin\frac1t\,\dd t,
\]
其中在 $t=0$ 的函数值可任意补定。由于 $G(0)=0$，
\[
  \abs*{\frac{G(x)-G(0)}x}
  \leqslant\frac1{\abs x}\int_0^{\abs x}t\,\dd t
  =\frac{\abs x}{2}\to0.
\]
故
\[
  \boxed{G'(0)=0}.
\]
\end{xsolution}

\paragraph{练习题 1（2）}
\begin{xsolution}
把 $\sin t/t$ 在 $t=0$ 处补定义为 $1$，得到连续函数
\[
  \varphi(t)=
  \begin{cases}
    \dfrac{\sin t}{t},&t\ne0,\\
    1,&t=0.
  \end{cases}
\]
于是由变限积分求导公式
\[
  \frac{\dd}{\dd x}\int_{x^2}^{x^3}\frac{\sin t}{t}\,\dd t
  =3x^2\varphi(x^3)-2x\varphi(x^2).
\]
因此当 $x\ne0$ 时可写为
\[
  \boxed{\frac{3\sin x^3-2\sin x^2}{x}},
\]
而在 $x=0$ 时导数为 $0$。
\end{xsolution}

\paragraph{练习题 1（3）}
\begin{xsolution}
这是 $0/0$ 型。由 L'Hospital 法则，
\[
  \lim_{x\to0}\frac{\int_0^x(\arctan t)^2\,\dd t}{x^2}
  =\lim_{x\to0}\frac{(\arctan x)^2}{2x}=0,
\]
因为 $\arctan x\sim x$。故极限为
\[
  \boxed{0}.
\]
\end{xsolution}

\paragraph{练习题 1（4）}
\begin{xsolution}
记
\[
  A(x)=\int_0^x\ee^{t^2}\,\dd t,
  \qquad
  B(x)=\int_0^x\ee^{2t^2}\,\dd t.
\]
当 $x\to+\infty$ 时二者都趋于 $+\infty$，且所求比值为 $A^2/B$。一次用 L'Hospital 法则，
\[
  \lim_{x\to+\infty}\frac{A^2(x)}{B(x)}
  =\lim_{x\to+\infty}2A(x)\ee^{-x^2}.
\]
又
\[
  \lim_{x\to+\infty}\frac{A(x)}{\ee^{x^2}}
  =\lim_{x\to+\infty}\frac{\ee^{x^2}}{2x\ee^{x^2}}
  =0.
\]
故
\[
  \boxed{\lim_{x\to+\infty}\frac{\left(\int_0^x\ee^{t^2}\,\dd t\right)^2}
  {\int_0^x\ee^{2t^2}\,\dd t}=0}.
\]
\end{xsolution}

\paragraph{练习题 2}
\begin{xsolution}
任取 $x\in[a,b)$，令 $0<h\leqslant b-x$。题设给出
\[
  \abs*{\frac1h\int_x^{x+h}f(t)\,\dd t}\leqslant Mh^\eta.
\]
因为 $f$ 在 $x$ 连续，由积分第一中值定理或变限积分基本定理，
\[
  \lim_{h\to0+}\frac1h\int_x^{x+h}f(t)\,\dd t=f(x).
\]
而右边 $Mh^\eta\to0$，故 $f(x)=0$。端点 $b$ 再由连续性得到 $f(b)=0$。所以
\[
  f\equiv0.
\]
\end{xsolution}

\paragraph{练习题 3}
\begin{xsolution}
令
\[
  F(x)=\int_a^x f(t)\,\dd t.
\]
则 $F(a)=0$，且 $F'(x)=f(x)$。题设即
\[
  F'(x)\leqslant F(x).
\]
于是
\[
  \bigl(\ee^{-x}F(x)\bigr)'=\ee^{-x}(F'(x)-F(x))\leqslant0.
\]
所以对 $x\in[a,b]$，
\[
  \ee^{-x}F(x)\leqslant\ee^{-a}F(a)=0,
\]
即 $F(x)\leqslant0$。最后由题设 $f(x)\leqslant F(x)$，得到
\[
  \boxed{f(x)\leqslant0\quad(x\in[a,b])}.
\]
\end{xsolution}

\paragraph{练习题 4}
\begin{xsolution}
先证明 $f$ 至少有一个零点。若 $f$ 在 $(0,\pi)$ 无零点，由连续性它在 $(0,\pi)$ 恒正或恒负。由于 $\sin t>0$，便有
\[
  \int_0^\pi f(t)\sin t\,\dd t\ne0,
\]
与题设矛盾。

反设 $f$ 在 $(0,\pi)$ 中恰有一个零点 $c$。由上面的正弦积分等于 $0$ 可知，$f$ 必须在 $c$ 两侧变号。函数
\[
  \varphi(t)=\sin(t-c)=\sin t\cos c-\cos t\sin c
\]
在 $(0,c)$ 与 $(c,\pi)$ 两侧也恰好异号。必要时将 $\varphi$ 乘以 $-1$，可使 $f(t)\varphi(t)\geqslant0$，并且除零点外严格大于 $0$。因此
\[
  \int_0^\pi f(t)\varphi(t)\,\dd t>0.
\]
但由题设两个正交条件，
\[
  \int_0^\pi f(t)\varphi(t)\,\dd t
  =\cos c\int_0^\pi f(t)\sin t\,\dd t
  -\sin c\int_0^\pi f(t)\cos t\,\dd t=0,
\]
矛盾。故 $f$ 在 $(0,\pi)$ 内至少有两个零点。
\end{xsolution}

\paragraph{练习题 5}
\begin{xsolution}
设 $T>0$ 是 $f$ 的一个周期，令
\[
  c=\frac1T\int_0^T f(t)\,\dd t,
  \qquad
  G(x)=\int_0^x(f(t)-c)\,\dd t.
\]
周期函数在任意长度为 $T$ 的区间上的积分相同，所以
\[
  \int_x^{x+T}(f(t)-c)\,\dd t
  =\int_0^T f(t)\,\dd t-cT=0.
\]
于是
\[
  G(x+T)-G(x)=0,
\]
即 $G$ 是周期为 $T$ 的周期函数。另一方面
\[
  F(x)=\int_0^x f(t)\,\dd t=G(x)+cx.
\]
故 $F$ 是一个周期函数与一个线性函数之和。
\end{xsolution}

\paragraph{练习题 6}
\begin{xsolution}
设
\[
  H(a)=\int_a^{a+T}f(x)\,\dd x.
\]
题设说明 $H$ 为常数。由于 $f$ 连续，可对变动积分限求导：
\[
  H'(a)=f(a+T)-f(a)=0.
\]
故对一切 $a\in\RR$，
\[
  f(a+T)=f(a),
\]
所以 $f$ 是周期为 $T$ 的周期函数。
\end{xsolution}

\paragraph{练习题 7}
\begin{xsolution}
对每个固定的 $b>0$，令
\[
  H_b(a)=\int_a^{ab}f(x)\,\dd x.
\]
题设说明 $H_b(a)$ 与 $a$ 无关，所以
\[
  0=H_b'(a)=b f(ab)-f(a).
\]
于是对任意 $a,b>0$，
\[
  b f(ab)=f(a).
\]
取 $a=1$，得到
\[
  f(b)=\frac{f(1)}b.
\]
因此全部解为
\[
  \boxed{f(x)=\frac Cx\quad(x>0)},
\]
其中 $C$ 为任意常数。反过来，
\[
  \int_a^{ab}\frac Cx\,\dd x=C\ln b,
\]
确与 $a$ 无关。
\end{xsolution}

\paragraph{练习题 8}
\begin{xsolution}
令
\[
  H(\xi)=\int_a^\xi f(x)\,\dd x-\int_\xi^b f(x)\,\dd x
  =2\int_a^\xi f(x)\,\dd x-\int_a^b f(x)\,\dd x.
\]
由命题 10.3.1，$H$ 连续。记 $I=\int_a^b f$，则
\[
  H(a)=-I,
  \qquad
  H(b)=I.
\]
由介值定理，存在 $\xi\in[a,b]$ 使 $H(\xi)=0$，即
\[
  \int_\xi^b f=\int_a^\xi f.
\]

这样的 $\xi$ 未必能取在开区间内。例如在 $[0,1]$ 上取
\[
  f(x)=1-2x.
\]
此时 $\int_0^1 f=0$，并且
\[
  \int_0^\xi f(x)\,\dd x=\xi-\xi^2>0
  \qquad(0<\xi<1).
\]
题中等式在总积分为 $0$ 时等价于 $\int_0^\xi f=0$，因此只有 $\xi=0,1$ 两个端点解，没有内点解。
\end{xsolution}

\paragraph{练习题 9}
\begin{xsolution}
由 $f>0$，有
\begin{align*}
  \frac{\dd}{\dd x}
  \left(\int_a^x f(t)\,\dd t-\int_x^b\frac{\dd t}{f(t)}\right)
  &=f(x)+\frac1{f(x)}\\
  &\geqslant2,
\end{align*}
其中最后一步是正数的算术--几何平均不等式。因此题中不等式成立。
\end{xsolution}

\paragraph{练习题 10}
\begin{xsolution}
记
\[
  A(x)=\int_1^x a(t)c(t)\,\dd t,\quad
  B(x)=\int_1^x b(t)d(t)\,\dd t,
\]
\[
  C(x)=\int_1^x a(t)d(t)\,\dd t,\quad
  D(x)=\int_1^x b(t)c(t)\,\dd t,
\]
并令 $E(x)=A(x)B(x)-C(x)D(x)$。由于 $a,b,c,d$ 都是多项式，$A,B,C,D,E$ 也都是多项式。

令 $h=x-1$，并简记
\[
  p=ac,\quad q=bd,\quad r=ad,\quad s=bc.
\]
在 $t=1$ 附近作 Taylor 展开，有
\[
  A(x)=p(1)h+\frac{p'(1)}2h^2+\bigO(h^3),
\]
其余三个积分同理。因此
\begin{align*}
  E(x)
  &=[p(1)q(1)-r(1)s(1)]h^2\\
  &\quad+\frac12[p'(1)q(1)+p(1)q'(1)-r'(1)s(1)-r(1)s'(1)]h^3
  +\bigO(h^4).
\end{align*}
第一项系数为
\[
  (ac)(bd)-(ad)(bc)=0.
\]
第二项括号展开为
\begin{align*}
 &(ac)'bd+ac(bd)'-(ad)'bc-ad(bc)'\\
 &=a'cbd+ac'bd+acb'd+acbd'\\
 &\quad-a'dbc-ad'bc-adb'c-adbc'=0,
\end{align*}
其中各量均取在 $t=1$，各项两两抵消。故
\[
  E(x)=\bigO((x-1)^4)\qquad(x\to1).
\]
而 $E$ 是多项式，所以 $x=1$ 是 $E$ 的至少四重零点，即
\[
  \boxed{(x-1)^4\mid E(x)}.
\]
\end{xsolution}

\paragraph{练习题 11}
\begin{xsolution}
令
\[
  F(x)=\int_0^x f(t)\,\dd t.
\]
则 $F'=f$，并且
\[
  g(x)=f(x)F(x)=\frac12(F^2(x))'.
\]
题设说 $g$ 单调减少，故 $(F^2)'$ 单调减少，于是 $F^2$ 是 $\RR$ 上的凹函数。

任何在整个实轴上非负的凹函数都只能是常数。事实上，若某两点 $x_1<x_2$ 有 $F^2(x_2)>F^2(x_1)$，则割线斜率为正；由凹性，向左延长时函数值最终会小于 $0$，矛盾。若 $F^2(x_2)<F^2(x_1)$，则向右延长同样得到矛盾。因此 $F^2$ 为常数。

又 $F(0)=0$，所以 $F^2\equiv0$，即 $F\equiv0$。于是
\[
  f=F'\equiv0.
\]
\end{xsolution}

\paragraph{练习题 12}
\begin{xsolution}
令
\[
  F(x)=\int_0^x f(t)\,\dd t.
\]
题设为
\[
  \int_0^x t f(t)\,\dd t=\frac x3F(x).
\]
对 $x>0$ 求导：
\[
  x f(x)=\frac13F(x)+\frac x3 f(x),
\]
即
\[
  F(x)=2x f(x).
\]
再求导，利用 $F'=f$，得
\[
  f(x)=2f(x)+2x f'(x),
\]
故
\[
  2x f'(x)+f(x)=0.
\]
在 $(0,+\infty)$ 上其一般解为
\[
  f(x)=\frac C{\sqrt x}.
\]
但题设 $f$ 在 $[0,+\infty)$ 上可微，特别在 $0$ 连续，因此只有 $C=0$ 才可能。故
\[
  \boxed{f\equiv0}.
\]
\end{xsolution}

\subsection*{10.4.6 练习题}

\paragraph{练习题 1}
\begin{xsolution}
错误在于
\[
  \frac1{\sqrt3}\arctan\left(\frac{\tan x}{\sqrt3}\right)
\]
虽在不越过 $\pi/2$ 的区间内是原函数，但它在 $x=\pi/2$ 处发生跳跃，不能在整个 $[0,3\pi/4]$ 上直接使用 Newton--Leibniz 公式。

直接令 $u=\cos x$，则 $\dd u=-\sin x\,\dd x$，于是
\begin{align*}
  \int_0^{3\pi/4}\frac{\sin x}{1+\cos^2x}\,\dd x
  &=\int_{-\sqrt2/2}^{1}\frac{\dd u}{1+u^2}\\
  &=\arctan1-\arctan\left(-\frac1{\sqrt2}\right)\\
  &=\boxed{\frac\pi4+\arctan\frac1{\sqrt2}}.
\end{align*}
\end{xsolution}

\paragraph{练习题 2（1）}
\begin{xsolution}
在 $x=1$ 处分段：
\[
  \int_0^2\abs{1-x}\,\dd x
  =\int_0^1(1-x)\,\dd x+\int_1^2(x-1)\,\dd x
  =\boxed{1}.
\]
\end{xsolution}

\paragraph{练习题 2（2）}
\begin{xsolution}
被积函数为偶函数，而且在 $x>0$ 时
\[
  x^2=\frac1x\Longleftrightarrow x=1.
\]
因此
\begin{align*}
  \int_{-2}^2\min\left\{\frac1{\abs x},x^2\right\}\dd x
  &=2\left(\int_0^1x^2\,\dd x+\int_1^2\frac{\dd x}{x}\right)\\
  &=\boxed{\frac23+2\ln2}.
\end{align*}
$x=0$ 处可按连续延拓取值 $0$，不影响积分。
\end{xsolution}

\paragraph{练习题 2（3）}
\begin{xsolution}
写成
\[
  \sqrt x\int_x^{x+1}\frac{\dd t}{\sqrt{t+\cos t}}
  =\int_x^{x+1}\sqrt{\frac{x}{t+\cos t}}\,\dd t.
\]
对 $t\in[x,x+1]$，有
\[
  \frac{t+\cos t}{x}=1+\bigO\left(\frac1x\right)
\]
且误差对该区间一致，因此被积函数一致趋于 $1$。积分区间长度为 $1$，故
\[
  \boxed{\lim_{x\to+\infty}\sqrt x\int_x^{x+1}
  \frac{\dd t}{\sqrt{t+\cos t}}=1}.
\]
\end{xsolution}

\paragraph{练习题 2（4）}
\begin{xsolution}
令
\[
  h(x)=\frac{\cos^2x}{1+\ee^{-x}}.
\]
则
\[
  h(x)+h(-x)
  =\cos^2x\left(\frac1{1+\ee^{-x}}+\frac1{1+\ee^x}\right)
  =\cos^2x.
\]
所以
\begin{align*}
  \int_{-\pi/4}^{\pi/4}h(x)\,\dd x
  &=\int_0^{\pi/4}\cos^2x\,\dd x\\
  &=\boxed{\frac\pi8+\frac14}.
\end{align*}
\end{xsolution}

\paragraph{练习题 2（5）}
\begin{xsolution}
令
\[
  q(t)=\frac{\sin t}{\pi-t},
\]
并在 $t=\pi$ 处补定义 $q(\pi)=1$，则 $q$ 连续。利用例题 10.3.1 的公式，
\begin{align*}
  \int_0^\pi\left(\int_0^x\frac{\sin t}{\pi-t}\,\dd t\right)\dd x
  &=\int_0^\pi q(x)(\pi-x)\,\dd x\\
  &=\int_0^\pi\sin x\,\dd x\\
  &=\boxed{2}.
\end{align*}
\end{xsolution}

\paragraph{练习题 2（6）}
\begin{xsolution}
被积函数关于 $\pi/2$ 对称，故
\[
  I=2\int_0^{\pi/2}\frac{\dd x}{a^2\sin^2x+b^2\cos^2x}.
\]
令 $t=\tan x$，则
\[
  I=2\int_0^{+\infty}\frac{\dd t}{a^2t^2+b^2}
  =\boxed{\frac\pi{\abs{ab}}}.
\]
这里 $ab\ne0$，与题设一致。
\end{xsolution}

\paragraph{练习题 3（1）}
\begin{xsolution}
令 $h(x)=1/(1+\cos^2x)$，则 $h(\pi-x)=h(x)$。由中点对称性，
\[
  \int_0^\pi xh(x)\,\dd x
  =\frac\pi2\int_0^\pi h(x)\,\dd x.
\]
又
\[
  \int_0^\pi h(x)\,\dd x
  =2\int_0^{\pi/2}\frac{\dd x}{1+\cos^2x}
  =2\int_0^{+\infty}\frac{\dd t}{t^2+2}
  =\frac\pi{\sqrt2},
\]
其中用了 $t=\tan x$。因此
\[
  \boxed{\int_0^\pi\frac{x\,\dd x}{1+\cos^2x}
  =\frac{\pi^2}{2\sqrt2}}.
\]
\end{xsolution}

\paragraph{练习题 3（2）}
\begin{xsolution}
令
\[
  h(x)=\frac1{\ee^x+\ee^{1-x}}.
\]
则 $h(1-x)=h(x)$，故
\[
  \int_0^1xh(x)\,\dd x=\frac12\int_0^1h(x)\,\dd x.
\]
令 $u=\ee^x$，则
\[
  \int_0^1h(x)\,\dd x
  =\int_1^{\ee}\frac{\dd u}{u^2+\ee}
  =\ee^{-1/2}\left(\arctan\sqrt{\ee}-\arctan\frac1{\sqrt\ee}\right).
\]
利用 $\arctan y+\arctan(1/y)=\pi/2$（$y>0$），得到
\[
  \boxed{\int_0^1\frac{x}{\ee^x+\ee^{1-x}}\,\dd x
  =\ee^{-1/2}\left(\arctan\sqrt{\ee}-\frac\pi4\right)}.
\]
\end{xsolution}

\paragraph{练习题 3（3）}
\begin{xsolution}
把 $x$ 与 $-x$ 配对：
\begin{align*}
  x\ln(1+\ee^x)+(-x)\ln(1+\ee^{-x})
  &=x\ln\frac{1+\ee^x}{1+\ee^{-x}}\\
  &=x\ln\ee^x=x^2.
\end{align*}
因此
\[
  \int_{-2}^2x\ln(1+\ee^x)\,\dd x
  =\int_0^2x^2\,\dd x
  =\boxed{\frac83}.
\]
\end{xsolution}

\paragraph{练习题 3（4）}
\begin{xsolution}
令 $I=\int_0^{\pi/4}\ln(1+\tan x)\,\dd x$。由 $x\mapsto\pi/4-x$，
\[
  \ln(1+\tan x)+
  \ln\left(1+\tan\left(\frac\pi4-x\right)\right)=\ln2.
\]
故
\[
  2I=\frac\pi4\ln2,
  \qquad
  \boxed{I=\frac\pi8\ln2}.
\]
\end{xsolution}

\paragraph{练习题 3（5）}
\begin{xsolution}
令
\[
  h(x)=\frac{\sin^n x}{\sin^n x+\cos^n x}.
\]
则
\[
  h\left(\frac\pi2-x\right)=1-h(x).
\]
故由对称性
\[
  \boxed{\int_0^{\pi/2}h(x)\,\dd x=\frac\pi4}.
\]
\end{xsolution}

\paragraph{练习题 3（6）}
\begin{xsolution}
令
\[
  A=a^n,\qquad B=b^n,
  \qquad p=\abs A,\quad q=\abs B.
\]
若 $A,B$ 不同时为 $0$，题中被积函数在可能的有限个 $0/0$ 点作可去延拓后可积。记
\[
  J_s=\int_0^\pi\frac{\sin^2x}{p^2\sin^2x+q^2\cos^2x}\,\dd x,
\]
\[
  J_c=\int_0^\pi\frac{\cos^2x}{p^2\sin^2x+q^2\cos^2x}\,\dd x.
\]
由
\[
  J_s+J_c=\frac\pi{pq},
  \qquad
  p^2J_s+q^2J_c=\pi
\]
（$pq>0$ 时），解得
\[
  J_s=\frac\pi{p(p+q)},
  \qquad
  J_c=\frac\pi{q(p+q)}.
\]
因此原积分为
\[
  AJ_s+BJ_c
  =\boxed{\frac{\pi(\operatorname{sgn}A+\operatorname{sgn}B)}{p+q}}.
\]
当 $a,b$ 同号，或 $n$ 为偶数时，这化为
\[
  \boxed{\frac{2\pi}{a^n+b^n}}.
\]
当 $ab<0$ 且 $n$ 为奇数时，积分为 $0$。若恰有 $a,b$ 中一个为 $0$，按可去延拓直接得到相应值 $\pi/a^n$ 或 $\pi/b^n$。若 $a=b=0$，原被积函数在整个区间上均无定义，因而原积分不存在。
\end{xsolution}

\paragraph{练习题 4（1）}
\begin{xsolution}
令
\[
  q(x)=\frac{f(x)}{f(x)+f(a-x)}.
\]
由题设分母不为 $0$，且
\[
  q(x)+q(a-x)=1.
\]
于是
\[
  2I=\int_0^a[q(x)+q(a-x)]\,\dd x=a,
\]
故
\[
  \boxed{I=\frac a2}.
\]
\end{xsolution}

\paragraph{练习题 4（2）}
\begin{xsolution}
在题中积分有定义，即 $1+f(x)\ne0$ 的前提下，令
\[
  q(x)=\frac1{1+f(x)}.
\]
由 $f(x)f(a-x)=1$，
\[
  q(a-x)=\frac1{1+f(a-x)}
  =\frac{f(x)}{1+f(x)}.
\]
所以 $q(x)+q(a-x)=1$，从而
\[
  \boxed{I=\frac a2}.
\]
若允许 $f(x)=-1$，则原被积函数可能无定义；例如 $f\equiv-1$。因此“积分有定义”是这一结论所需的隐含条件。
\end{xsolution}

\paragraph{练习题 5（1）}
\begin{xsolution}
令 $h(x)=f(\sin x)$，则 $h(\pi-x)=h(x)$。于是
\[
  \int_0^\pi xh(x)\,\dd x
  =\frac\pi2\int_0^\pi h(x)\,\dd x,
\]
即所证等式。
\end{xsolution}

\paragraph{练习题 5（2）}
\begin{xsolution}
原积分含有 $\dd x/x$，故按通常的实积分意义需有 $a>0$；以下在这个使题目有定义的条件下证明。左端令
\[
  u=\frac{x^2}{a},
  \qquad
  \frac{\dd x}{x}=\frac12\frac{\dd u}{u},
\]
得到
\[
  L=\frac12\int_{1/a}^{a}f\left(a\left(u+\frac1u\right)\right)\frac{\dd u}{u}.
\]
把区间拆成 $[1/a,1]$ 与 $[1,a]$，第二段作代换 $u=1/v$，可见两段相等，因此
\[
  L=\int_{1/a}^{1}f\left(a\left(u+\frac1u\right)\right)\frac{\dd u}{u}.
\]
右端令 $u=x/a$，同样得到
\[
  R=\int_{1/a}^{1}f\left(a\left(u+\frac1u\right)\right)\frac{\dd u}{u}.
\]
故 $L=R$。
\end{xsolution}

\paragraph{练习题 5（3）}
\begin{xsolution}
令
\[
  r=\sqrt{a^2+b^2}.
\]
可取 $\varphi$ 使
\[
  a\cos x+b\sin x=r\cos(x-\varphi).
\]
由 $2\pi$ 周期性作平移，
\[
  \int_0^{2\pi}f(a\cos x+b\sin x)\,\dd x
  =\int_0^{2\pi}f(r\cos u)\,\dd u.
\]
而 $u\mapsto2\pi-u$ 使后半区间的积分等于前半区间，故
\[
  \int_0^{2\pi}f(r\cos u)\,\dd u
  =2\int_0^\pi f(r\cos u)\,\dd u,
\]
即所证。
\end{xsolution}

\paragraph{练习题 6}
\begin{xsolution}
对第一个积分，被积函数
\[
  \sin^{2n-1}x\cos(2n+1)x
\]
在变换 $x\mapsto\pi-x$ 下变号，因为 $\sin(\pi-x)=\sin x$，而
\[
  \cos((2n+1)(\pi-x))=-\cos((2n+1)x).
\]
故第一个积分为 $0$。

对第二个积分，$\cos^{2n-1}(\pi-x)=-\cos^{2n-1}x$，而
\[
  \sin((2n+1)(\pi-x))=\sin((2n+1)x),
\]
所以第二个被积函数也关于 $\pi/2$ 反对称，积分同样为 $0$。因此
\[
  \boxed{0,\quad0}.
\]
\end{xsolution}

\paragraph{练习题 7}
\begin{xsolution}
令 $x=\ee^{-t}$，则 $\dd x=-\ee^{-t}\dd t$，并且 $x:0\to1$ 对应 $t:+\infty\to0$。于是
\begin{align*}
  I(m,n)
  &=(-1)^n\int_0^{+\infty}t^n\ee^{-(m+1)t}\,\dd t.
\end{align*}
连续分部积分 $n$ 次可得
\[
  \int_0^{+\infty}t^n\ee^{-\lambda t}\,\dd t
  =\frac{n!}{\lambda^{n+1}},\qquad\lambda>0.
\]
故
\[
  \boxed{I(m,n)=\frac{(-1)^n n!}{(m+1)^{n+1}}}.
\]
\end{xsolution}

\paragraph{练习题 8}
\begin{xsolution}
记
\[
  J(m,n)=\int_0^{\pi/2}\sin^m x\cos^n x\,\dd x.
\]
分部积分得到递推式
\[
  J(m,n)=\frac{m-1}{m+n}J(m-2,n)\qquad(m\geqslant2),
\]
以及对称关系 $J(m,n)=J(n,m)$。反复递推到 $0$ 或 $1$，利用
\[
  J(0,0)=\frac\pi2,
  \qquad
  J(1,0)=J(0,1)=1,
  \qquad
  J(1,1)=\frac12,
\]
得到
\[
  \boxed{
  J(m,n)=
  \begin{cases}
    \displaystyle
    \frac{(m-1)!!(n-1)!!}{(m+n)!!}\frac\pi2,
    &m,n\text{ 都为偶数},\\[8pt]
    \displaystyle
    \frac{(m-1)!!(n-1)!!}{(m+n)!!},
    &m,n\text{ 至少一个为奇数}.
  \end{cases}}
\]
其中约定 $(-1)!!=0!!=1$。
\end{xsolution}

\paragraph{练习题 9}
\begin{xsolution}
设
\[
  F(x)=\int_0^x\cos\frac1t\,\dd t,
  \qquad F(0)=0.
\]
对 $x\ne0$，利用
\[
  \dd\left(\sin\frac1t\right)=-\frac1{t^2}\cos\frac1t\,\dd t
\]
作分部积分，得
\[
  F(x)=-x^2\sin\frac1x+\int_0^x2t\sin\frac1t\,\dd t.
\]
因而
\[
  \abs{F(x)}\leqslant x^2+\abs{x}^2=2x^2.
\]
所以
\[
  F'(0)=\lim_{x\to0}\frac{F(x)}x=0.
\]
即
\[
  \boxed{F'(0)=0}.
\]
\end{xsolution}

\paragraph{练习题 10（Fejér 积分）}
\begin{xsolution}
利用有限 Fourier 恒等式
\[
  \left(\frac{\sin nx}{\sin x}\right)^2
  =n+2\sum_{k=1}^{n-1}(n-k)\cos2kx.
\]
该等式在 $x=0$ 处按连续延拓理解。积分 $0$ 到 $\pi/2$，对每个 $k\geqslant1$ 有
\[
  \int_0^{\pi/2}\cos2kx\,\dd x=0.
\]
因此
\[
  \boxed{\int_0^{\pi/2}\left(\frac{\sin nx}{\sin x}\right)^2\dd x
  =\frac{n\pi}{2}}.
\]
\end{xsolution}

\paragraph{练习题 11（1）}
\begin{xsolution}
因为 $\abs{\sin t}$ 以 $\pi$ 为周期，
\begin{align*}
  f(x+\pi)
  &=\int_{x+\pi}^{x+3\pi/2}\abs{\sin t}\,\dd t\\
  &=\int_x^{x+\pi/2}\abs{\sin(u+\pi)}\,\dd u
  =f(x).
\end{align*}
故 $f$ 以 $\pi$ 为周期。
\end{xsolution}

\paragraph{练习题 11（2）}
\begin{xsolution}
只需考察 $x\in[0,\pi]$。当 $0\leqslant x\leqslant\pi/2$ 时，
\[
  f(x)=\int_x^{x+\pi/2}\sin t\,\dd t
  =\sin x+\cos x.
\]
故在这一段最大值为 $\sqrt2$（$x=\pi/4$），端点值均为 $1$。

利用 $f'(x)=\abs{\cos x}-\abs{\sin x}$，在 $[\pi/2,\pi]$ 上唯一的内点极值出现在 $x=3\pi/4$，且为极小值。直接计算
\[
  f\left(\frac{3\pi}4\right)
  =\int_{3\pi/4}^{5\pi/4}\abs{\sin t}\,\dd t
  =2\int_{3\pi/4}^{\pi}\sin t\,\dd t
  =2-\sqrt2.
\]
因此
\[
  \boxed{\max f=\sqrt2,\qquad \min f=2-\sqrt2}.
\]
\end{xsolution}

\paragraph{练习题 12}
\begin{xsolution}
由积分第一中值定理，存在 $c\in(7/8,1)$，使
\[
  \int_{7/8}^1f(x)\,\dd x=\frac18f(c).
\]
与题设比较得 $f(c)=f(0)$。函数 $f$ 在 $[0,c]$ 连续、在 $(0,c)$ 可微，由 Rolle 定理，存在 $\xi\in(0,c)\subset(0,1)$，使
\[
  \boxed{f'(\xi)=0}.
\]
\end{xsolution}


\subsection*{10.5.2 第一组参考题}

\paragraph{第一组参考题 1}
\begin{xsolution}
按定义取等距分割
\[
  x_i=a+i\Delta,
  \qquad
  \Delta=\frac{b-a}{n},
  \qquad i=0,1,\ldots,n,
\]
并取右端点为介点。对应 Riemann 和为
\[
  S_n=\Delta\sum_{i=1}^n(a+i\Delta)^m.
\]
由二项式公式
\[
  S_n=\sum_{k=0}^m C_m^k a^{m-k}\Delta^{k+1}
  \sum_{i=1}^n i^k.
\]
利用整数幂和的首项渐近式
\[
  \sum_{i=1}^n i^k=\frac{n^{k+1}}{k+1}+\bigO(n^k),
\]
且 $\Delta=(b-a)/n$，可得
\[
  \lim_{n\to\infty}S_n
  =\sum_{k=0}^m C_m^k a^{m-k}
  \frac{(b-a)^{k+1}}{k+1}.
\]
另一方面，将 $(a+y)^m$ 展开并从 $0$ 到 $b-a$ 积分，或直接使用二项式恒等式，右端即为
\[
  \frac{b^{m+1}-a^{m+1}}{m+1}.
\]
由于 $x^m$ 连续，Riemann 和的极限就是定积分。因此
\[
  \boxed{\int_a^b x^m\,\dd x=
  \frac{b^{m+1}-a^{m+1}}{m+1}}.
\]
\end{xsolution}

\paragraph{第一组参考题 2（1）}
\begin{xsolution}
在任意非退化区间 $[a,b]$ 上令
\[
  f(x)=
  \begin{cases}
    1,&x\in\QQ,\\
    -1,&x\notin\QQ.
  \end{cases}
\]
则 $f$ 处处不连续，故 $f\notin R[a,b]$，而
\[
  \abs{f(x)}\equiv1\in R[a,b].
\]
这说明 $\abs f$ 可积不能推出 $f$ 可积。
\end{xsolution}

\paragraph{第一组参考题 2（2）}
\begin{xsolution}
设 $f$ 是某个函数 $F$ 的导函数，并且 $\abs f\in R[a,b]$。于是 $\abs f$ 有界，从而 $f$ 也有界。只需证明 $f$ 的不连续点集包含在 $\abs f$ 的不连续点集中。

设 $x_0$ 是 $\abs f$ 的连续点。若 $f(x_0)=0$，则
\[
  \abs{f(x)-f(x_0)}=\abs{f(x)}\to0,
\]
所以 $f$ 在 $x_0$ 连续。

若 $f(x_0)\ne0$，则由 $\abs f$ 的连续性，存在邻域 $U$ 和 $c>0$，使
\[
  \abs{f(x)}\geqslant c\qquad(x\in U).
\]
导函数具有 Darboux 性质。若 $f$ 在 $U$ 内取过正、负两种符号，Darboux 性质会迫使它在中间取值 $0$，与上式矛盾。因此 $f$ 在足够小的 $U$ 上符号固定，于是 $f=\pm\abs f$，故 $f$ 在 $x_0$ 连续。

所以
\[
  D(f)\subseteq D(\abs f).
\]
右边为零测度集，$f$ 又有界，由 Lebesgue 判别定理得
\[
  \boxed{f\in R[a,b]}.
\]
\end{xsolution}

\paragraph{第一组参考题 3}
\begin{xsolution}
记
\[
  S_P=\sum_{k=1}^n f(\xi_k)g(\xi'_k)\Delta x_k,
  \qquad
  T_P=\sum_{k=1}^n f(\xi_k)g(\xi_k)\Delta x_k.
\]
设 $\abs{f}\leqslant M$。在第 $k$ 个子区间上记 $g$ 的振幅为 $\omega_k(g)$，则
\begin{align*}
  \abs{S_P-T_P}
  &\leqslant M\sum_{k=1}^n
  \abs{g(\xi'_k)-g(\xi_k)}\Delta x_k\\
  &\leqslant M\sum_{k=1}^n\omega_k(g)\Delta x_k.
\end{align*}
因为 $g\in R[a,b]$，由可积的第一充分必要条件，右端在 $\norm P\to0$ 时趋于 $0$。另一方面 $fg\in R[a,b]$，所以
\[
  T_P\to\int_a^b f(x)g(x)\,\dd x.
\]
因此
\[
  \boxed{\lim_{\norm P\to0}S_P=
  \int_a^b f(x)g(x)\,\dd x}.
\]
\end{xsolution}

\paragraph{第一组参考题 4}
\begin{xsolution}
设 $f$ 在 $(a,b)$ 上下凸。凸函数的基本性质给出：对每个 $x\in(a,b)$，左右导数 $f'_-(x),f'_+(x)$ 都存在，而且
\[
  f'_-(x)\leqslant f'_+(x),
\]
并且这两个函数都关于 $x$ 单调增加。在任意固定的 $[c,d]\subset(a,b)$ 上，它们因此有界且单调，故均 Riemann 可积。

取分割
\[
  c=x_0<x_1<\cdots<x_n=d.
\]
由凸函数割线斜率的单调性，对每个 $i$ 有
\[
  f'_+(x_{i-1})
  \leqslant
  \frac{f(x_i)-f(x_{i-1})}{x_i-x_{i-1}}
  \leqslant f'_-(x_i).
\]
乘以 $\Delta x_i$ 并求和：
\[
  \sum_{i=1}^n f'_+(x_{i-1})\Delta x_i
  \leqslant f(d)-f(c)
  \leqslant
  \sum_{i=1}^n f'_-(x_i)\Delta x_i.
\]
令 $\norm P\to0$，左右两端分别趋于相应积分，得到
\[
  \int_c^d f'_+(x)\,\dd x
  \leqslant f(d)-f(c)
  \leqslant
  \int_c^d f'_-(x)\,\dd x.
\]
又因逐点 $f'_-(x)\leqslant f'_+(x)$，
\[
  \int_c^d f'_-(x)\,\dd x
  \leqslant\int_c^d f'_+(x)\,\dd x.
\]
故三者必相等，即
\[
  \boxed{f(d)-f(c)=\int_c^d f'_-(x)\,\dd x
  =\int_c^d f'_+(x)\,\dd x}.
\]
\end{xsolution}

\paragraph{第一组参考题 5（1）}
\begin{xsolution}
设 $f\in R[a,b]$，给定 $\varepsilon>0$。由可积的第二充分必要条件，可取分割
\[
  P:\ a=x_0<\cdots<x_n=b
\]
使
\[
  \sum_{i=1}^n\omega_i\Delta x_i<\varepsilon,
\]
其中 $\omega_i$ 是 $f$ 在第 $i$ 个子区间上的振幅。在每个子区间内任取 $\xi_i$，定义阶梯函数
\[
  g(x)=f(\xi_i),\qquad x\in(x_{i-1},x_i),
\]
分点处任意补定义。则
\[
  \abs{f(x)-g(x)}\leqslant\omega_i
\]
在第 $i$ 个子区间成立，所以
\[
  \int_a^b\abs{f-g}\leqslant
  \sum_{i=1}^n\omega_i\Delta x_i<\varepsilon.
\]
\end{xsolution}

\paragraph{第一组参考题 5（2）--（3）}
\begin{xsolution}
从（1）的阶梯函数 $g$ 出发。设它只有有限个跳跃点，且 $\abs g\leqslant M$。在每个跳跃点两侧取互不相交的小区间，使这些小区间总长度小于 $\delta$；在区间外保持 $g$ 不变，在每个小区间内用直线连接左右常值，得到连续折线函数 $h$。则 $g-h$ 只在这些小区间内非零，并且有界，因此
\[
  \int_a^b\abs{g-h}\,\dd x\leqslant 2M\delta.
\]
先在（1）中把误差取为 $\varepsilon/2$，再取 $\delta$ 使 $2M\delta<\varepsilon/2$，便有
\[
  \int_a^b\abs{f-h}\,\dd x<\varepsilon.
\]
同一个 $h$ 既是折线函数，也是连续函数，因此同时解决（2）、（3）。
\end{xsolution}

\paragraph{第一组参考题 5（4）}
\begin{xsolution}
若 $a=b$，结论平凡。以下设 $a<b$。继续用（2）--（3）所得的连续函数 $h$，并先把前一步误差取得满足
\[
  \int_a^b\abs{f(x)-h(x)}\,\dd x<\frac\varepsilon2.
\]
由 Weierstrass 多项式逼近定理，可取多项式 $q$，使
\[
  \max_{x\in[a,b]}\abs{h(x)-q(x)}
  <\frac{\varepsilon}{2(b-a)}.
\]
于是
\[
  \int_a^b\abs{h(x)-q(x)}\,\dd x<\frac\varepsilon2.
\]
因此
\[
  \int_a^b\abs{f(x)-q(x)}\,\dd x<\varepsilon.
\]
多项式 $q$ 当然属于 $C^1[a,b]$，故（4）得证。
\end{xsolution}

\paragraph{第一组参考题 6（积分的连续性命题）}
\begin{xsolution}
由第一组参考题 5（1），在稍大的区间 $[a-\delta,b+\delta]$ 上，对给定 $\varepsilon>0$ 可取阶梯函数 $g$，使
\[
  \int_{a-\delta}^{b+\delta}\abs{f(x)-g(x)}\,\dd x<\varepsilon.
\]
当 $\abs h<\delta$ 时，由三角不等式
\begin{align*}
  \int_a^b\abs{f(x+h)-f(x)}\,\dd x
  &\leqslant
  \int_a^b\abs{f(x+h)-g(x+h)}\,\dd x\\
  &\quad+
  \int_a^b\abs{g(x+h)-g(x)}\,\dd x
  +\int_a^b\abs{g(x)-f(x)}\,\dd x\\
  &\leqslant2\varepsilon+
  \int_a^b\abs{g(x+h)-g(x)}\,\dd x.
\end{align*}
阶梯函数 $g$ 只有有限个跳跃点。若 $\abs g\leqslant M$，则 $g(x+h)\ne g(x)$ 只能发生在这些跳跃点的 $\bigO(\abs h)$ 邻域内，故存在常数 $C$ 使
\[
  \int_a^b\abs{g(x+h)-g(x)}\,\dd x\leqslant C\abs h\to0.
\]
于是
\[
  \limsup_{h\to0}\int_a^b\abs{f(x+h)-f(x)}\,\dd x
  \leqslant2\varepsilon.
\]
由 $\varepsilon$ 任意，得到所证极限为 $0$。
\end{xsolution}

\paragraph{第一组参考题 7}
\begin{xsolution}
给定 $\varepsilon>0$。由可积的第二充分必要条件，可取分割 $P$，使
\[
  \sum_{i=1}^n\omega_i\Delta x_i<\varepsilon(b-a).
\]
若每个子区间的振幅都满足 $\omega_i\geqslant\varepsilon$，则
\[
  \sum_{i=1}^n\omega_i\Delta x_i
  \geqslant\varepsilon\sum_{i=1}^n\Delta x_i
  =\varepsilon(b-a),
\]
矛盾。因此至少有一个子区间 $[c,d]$ 满足
\[
  \boxed{\omega_{f|[c,d]}<\varepsilon}.
\]
\end{xsolution}

\paragraph{第一组参考题 8}
\begin{xsolution}
由 $f\in R[a,b]$ 和 Lebesgue 判别定理，$f$ 的不连续点集 $D(f)$ 为零测度集。任意非空开子区间 $(c,d)\subset[a,b]$ 的长度 $d-c>0$，不可能完全包含在零测度集 $D(f)$ 中。因此 $(c,d)$ 中必有一点不属于 $D(f)$，即必有 $f$ 的连续点。故连续点在 $[a,b]$ 中稠密。
\end{xsolution}

\paragraph{第一组参考题 9}
\begin{xsolution}
先证必要性。设 $f\geqslant0$ 且
\[
  \int_a^b f(x)\,\dd x=0.
\]
若 $x_0$ 是连续点且 $f(x_0)>0$，则存在小区间 $[c,d]$ 和 $\mu>0$，使 $f(x)\geqslant\mu$ 于 $[c,d]$，从而
\[
  \int_a^b f\geqslant\int_c^d f\geqslant\mu(d-c)>0,
\]
矛盾。所以每个连续点都满足 $f(x_0)=0$。

再证充分性。若 $f$ 在所有连续点取值 $0$，由于 $f$ Riemann 可积，其不连续点集为零测度集，所以 $f=0$ 几乎处处。设 $\abs f\leqslant M$。给定 $\varepsilon>0$，用总长度小于 $\varepsilon/M$ 的开区间族覆盖全部不连续点；在其余点 $f=0$。从而积分的绝对值不超过 $\varepsilon$。令 $\varepsilon\to0$，得
\[
  \int_a^b f=0.
\]
故条件充分必要。
\end{xsolution}

\paragraph{第一组参考题 10}
\begin{xsolution}
令 $h=f-g$。则 $h\in R[a,b]$，并且题设变为：在每个子区间中都存在一点使 $h(x)=0$。

若 $x_0$ 是 $h$ 的连续点而 $h(x_0)>0$，则由连续性存在包含 $x_0$ 的小区间，其中 $h>0$，这与该小区间中应有零点矛盾。同理不可能 $h(x_0)<0$。因此 $h$ 在所有连续点上都为 $0$。

由第一组参考题 9 的充分性证明（不要求非负时可对 $\abs h$ 使用同样的覆盖估计），有
\[
  \int_a^b h(x)\,\dd x=0.
\]
所以
\[
  \boxed{\int_a^b f(x)\,\dd x=\int_a^b g(x)\,\dd x}.
\]
\end{xsolution}

\paragraph{第一组参考题 11（积分第一中值定理的一种推广）}
\begin{xsolution}
不妨设 $g\geqslant0$；若 $g\leqslant0$，把 $g$ 换成 $-g$ 即可。若 $\int_a^b g=0$，因 $g\geqslant0$ 且可积，$fg$ 的积分也为 $0$，任取 $\xi\in(a,b)$ 即可。以下设
\[
  G:=\int_a^b g(x)\,\dd x>0.
\]
由积分第一中值定理，存在
\[
  \eta\in[m,M],
  \qquad
  m=\inf_{[a,b]}f,\quad M=\sup_{[a,b]}f,
\]
使
\[
  \int_a^b f(x)g(x)\,\dd x=\eta G.
\]
只需证明存在内点 $\xi$ 使 $f(\xi)=\eta$。

因为 $f$ 有原函数，所以 $f$ 是导函数，具有 Darboux 性质。若 $m<\eta<M$，由下确界、上确界定义可找到两点使函数值分别小于、 大于 $\eta$，再由 Darboux 性质可在二者之间找到内点 $\xi$ 使 $f(\xi)=\eta$。

若 $\eta=m<M$，则
\[
  \int_a^b(f(x)-m)g(x)\,\dd x=0,
\]
而两个因子都非负。由 $G>0$ 和例题 10.2.1，存在子区间 $[c,d]\subset(a,b)$ 与 $\mu>0$，使
\[
  g(x)\geqslant\mu\qquad(x\in[c,d]).
\]
于是
\[
  0\geqslant\mu\int_c^d(f(x)-m)\,\dd x\geqslant0,
\]
所以 $\int_c^d(f-m)=0$。函数 $f$ Riemann 可积，故在 $(c,d)$ 中有连续点 $\xi$；由非负函数积分为零的必要性，$f(\xi)-m=0$。因此 $f(\xi)=m=\eta$。$\eta=M$ 的情形同理。若 $m=M$，$f$ 为常数，也显然成立。

故总存在 $\xi\in(a,b)$ 使
\[
  \boxed{\int_a^b f(x)g(x)\,\dd x
  =f(\xi)\int_a^b g(x)\,\dd x}.
\]
\end{xsolution}

\paragraph{第一组参考题 12}
\begin{xsolution}
在 $x=1$ 附近作二阶展开：
\[
  \frac1{1+x}
  =\frac12-\frac14(x-1)+\bigO((x-1)^2).
\]
因此
\begin{align*}
  \int_0^1\frac{x^{n-1}}{1+x}\,\dd x
  &=\frac12\int_0^1x^{n-1}\,\dd x
  -\frac14\int_0^1x^{n-1}(x-1)\,\dd x\\
  &\quad+\bigO\left(\int_0^1x^{n-1}(1-x)^2\,\dd x\right).
\end{align*}
前两项为
\[
  \frac1{2n}
  +\frac1{4n(n+1)}.
\]
而
\[
  \int_0^1x^{n-1}(1-x)^2\,\dd x
  =\frac{2}{n(n+1)(n+2)}=\bigO(n^{-3}).
\]
所以
\[
  \int_0^1\frac{x^{n-1}}{1+x}\,\dd x
  =\frac1{2n}+\frac1{4n^2}+\bigO(n^{-3}).
\]
故
\[
  \boxed{a=\frac12,\qquad b=\frac14}.
\]
\end{xsolution}

\paragraph{第一组参考题 13}
\begin{xsolution}
由题设
\[
  f(x_p)=
  \left(\frac1{b-a}\int_a^b f^p(t)\,\dd t\right)^{1/p}.
\]
例题 10.2.5 给出
\[
  \left(\int_a^b f^p(t)\,\dd t\right)^{1/p}\to f(b),
\]
而 $(b-a)^{-1/p}\to1$，因此
\[
  f(x_p)\to f(b).
\]
$f$ 在 $[a,b]$ 上连续且严格单调增加，所以其反函数在值域上连续。于是
\[
  \boxed{\lim_{p\to+\infty}x_p=b}.
\]
\end{xsolution}

\paragraph{第一组参考题 14}
\begin{xsolution}
令
\[
  F(x)=\int_0^x f(t)\,\dd t,
\]
并设
\[
  f(x)+aF(x)\to L.
\]
则 $F'=f$，所以
\[
  F'(x)+aF(x)=L+r(x),
  \qquad r(x)\to0.
\]
令
\[
  y(x)=F(x)-\frac La.
\]
便有
\[
  y'(x)+ay(x)=r(x).
\]
对任意固定 $X$，乘积分因子后得到
\[
  y(x)=\ee^{-a(x-X)}y(X)
  +\int_X^x\ee^{-a(x-t)}r(t)\,\dd t.
\]
给定 $\varepsilon>0$，取 $X$ 足够大，使 $\abs{r(t)}<\varepsilon$（$t\geqslant X$）。于是
\[
  \limsup_{x\to\infty}\abs{y(x)}\leqslant\frac\varepsilon a.
\]
由 $\varepsilon$ 任意，$y(x)\to0$，即
\[
  F(x)\to\frac La.
\]
最后
\[
  f(x)=[f(x)+aF(x)]-aF(x)\to L-L=0.
\]
故
\[
  \boxed{f(+\infty)=0}.
\]
\end{xsolution}

\paragraph{第一组参考题 15（1）--（2）}
\begin{xsolution}
先作代换 $u=x+t$：
\[
  F(x)=\int_{x+a}^{x+b}f(u)\cos(u-x)\,\dd u.
\]
这里关于 $x$ 的被积表达式
\[
  \Phi(u,x)=f(u)\cos(u-x)
\]
连续，且
\[
  \frac{\partial\Phi}{\partial x}=f(u)\sin(u-x)
\]
也连续。因此可用含变动积分限的求导公式，得到 $F$ 在 $[a,b]$ 上可导，并且
\begin{align*}
  F'(x)
  &=f(x+b)\cos b-f(x+a)\cos a
  +\int_{x+a}^{x+b}f(u)\sin(u-x)\,\dd u\\
  &=\boxed{f(x+b)\cos b-f(x+a)\cos a
  +\int_a^b f(x+t)\sin t\,\dd t}.
\end{align*}
这同时完成了可导性的证明与导数计算。
\end{xsolution}

\paragraph{第一组参考题 16}
\begin{xsolution}
利用有限三角和恒等式。

若 $n=2m+1$ 为奇数，则
\[
  \frac{\sin(2m+1)x}{\sin x}
  =1+2\sum_{k=1}^m\cos2kx.
\]
在 $[0,\pi/2]$ 上积分，所有余弦项积分都为 $0$，故
\[
  \boxed{\int_0^{\pi/2}\frac{\sin nx}{\sin x}\,\dd x=\frac\pi2
  \quad(n\text{ 为奇数})}.
\]

若 $n=2m$ 为偶数，则
\[
  \frac{\sin2mx}{\sin x}
  =2\sum_{k=1}^m\cos(2k-1)x.
\]
因此
\[
  \boxed{\int_0^{\pi/2}\frac{\sin nx}{\sin x}\,\dd x
  =2\sum_{k=1}^{m}\frac{(-1)^{k-1}}{2k-1}
  \quad(n=2m)}.
\]
\end{xsolution}

\paragraph{第一组参考题 17（1）--（2）}
\begin{xsolution}
由二项式公式
\begin{align*}
  B(m,n)
  &=\sum_{k=0}^nC_n^k\frac{(-1)^k}{m+k+1}\\
  &=\sum_{k=0}^nC_n^k(-1)^k\int_0^1x^{m+k}\,\dd x\\
  &=\int_0^1x^m(1-x)^n\,\dd x.
\end{align*}
作代换 $x=1-t$，立即得到
\[
  B(m,n)=B(n,m),
\]
这证明了（1）。

连续分部积分，或使用例题 10.4.10 的计算，得到
\[
  \int_0^1x^m(1-x)^n\,\dd x
  =\frac{m!n!}{(m+n+1)!}.
\]
因此
\[
  \boxed{B(m,n)=B(n,m)=\frac{m!n!}{(m+n+1)!}}.
\]
\end{xsolution}

\paragraph{第一组参考题 18}
\begin{xsolution}
利用
\[
  \dd\left(\cos\frac1t\right)
  =\frac1{t^2}\sin\frac1t\,\dd t,
\]
对 $x>0$ 分部积分：
\[
  \int_0^x\sin\frac1t\,\dd t
  =x^2\cos\frac1x-
  \int_0^x2t\cos\frac1t\,\dd t.
\]
故
\[
  \abs*{\int_0^x\sin\frac1t\,\dd t}
  \leqslant x^2+x^2=2x^2.
\]
当 $m<2$ 时
\[
  \abs*{\frac1{x^m}\int_0^x\sin\frac1t\,\dd t}
  \leqslant2x^{2-m}\to0.
\]
所以所求极限为
\[
  \boxed{0}.
\]
\end{xsolution}

\paragraph{第一组参考题 19}
\begin{xsolution}
在 $0\leqslant\theta\leqslant\pi/2$ 上有弦线估计
\[
  \sin\theta\geqslant\frac{2\theta}{\pi}.
\]
因此
\begin{align*}
  0\leqslant
  \int_0^{\pi/2}\ee^{-R\sin\theta}\,\dd\theta
  &\leqslant\int_0^{\pi/2}\ee^{-2R\theta/\pi}\,\dd\theta\\
  &\leqslant\frac\pi{2R}.
\end{align*}
于是当 $\lambda<1$ 时
\[
  0\leqslant R^\lambda\int_0^{\pi/2}\ee^{-R\sin\theta}\,\dd\theta
  \leqslant\frac\pi2R^{\lambda-1}\to0.
\]
故极限为 $\boxed{0}$。
\end{xsolution}

\paragraph{第一组参考题 20}
\begin{xsolution}
记
\[
  I=\int_0^\pi(f(x)+f''(x))\sin x\,\dd x.
\]
对含 $f''$ 的积分分部积分两次：
\begin{align*}
  \int_0^\pi f''(x)\sin x\,\dd x
  &=-\int_0^\pi f'(x)\cos x\,\dd x\\
  &=-\bigl[f(x)\cos x\bigr]_0^\pi
  -\int_0^\pi f(x)\sin x\,\dd x.
\end{align*}
所以
\[
  I=-\bigl[f(x)\cos x\bigr]_0^\pi
  =f(\pi)+f(0).
\]
题设 $I=5$、$f(\pi)=2$，故
\[
  \boxed{f(0)=3}.
\]
\end{xsolution}

\paragraph{第一组参考题 21}
\begin{xsolution}
不存在满足条件的非负连续函数。

若存在，则由三个矩条件
\begin{align*}
  \int_0^1(x-a)^2f(x)\,\dd x
  &=\int_0^1x^2f(x)\,\dd x
  -2a\int_0^1xf(x)\,\dd x
  +a^2\int_0^1f(x)\,\dd x\\
  &=a^2-2a^2+a^2=0.
\end{align*}
被积函数 $(x-a)^2f(x)$ 连续且非负，所以处处为 $0$。若 $x\ne a$，便有 $f(x)=0$。由连续性，即使 $a\in[0,1]$，也只能有 $f(a)=0$。故 $f\equiv0$，与 $\int_0^1f=1$ 矛盾。因此对任何给定实数 $a$ 都无此非负连续函数。
\end{xsolution}

\paragraph{第一组参考题 22}
\begin{xsolution}
令
\[
  F(x)=\ee^x f(x).
\]
则
\[
  F'(x)=\ee^x(f(x)+f'(x)).
\]
题设等式乘以 $\ee$ 后成为
\[
  F(1)=3\int_0^{1/3}F(x)\,\dd x.
\]
由积分第一中值定理，存在 $c\in(0,1/3)$，使
\[
  3\int_0^{1/3}F(x)\,\dd x=F(c).
\]
故 $F(1)=F(c)$。由 Rolle 定理，存在 $\xi\in(c,1)\subset(0,1)$，使 $F'(\xi)=0$，从而
\[
  \boxed{f(\xi)+f'(\xi)=0}.
\]
\end{xsolution}

\paragraph{第一组参考题 23}
\begin{xsolution}
令
\[
  G(t)=\sqrt{1+2\int_0^t f(s)\,\dd s}.
\]
因 $f\geqslant0$，根号内至少为 $1$，且题设正好给出
\[
  f(t)\leqslant G(t).
\]
求导得
\[
  G'(t)=\frac{f(t)}{G(t)}\leqslant1.
\]
又 $G(0)=1$，所以
\[
  G(t)\leqslant1+t.
\]
结合 $f(t)\leqslant G(t)$，得到
\[
  \boxed{f(t)\leqslant1+t\qquad(0\leqslant t\leqslant1)}.
\]
\end{xsolution}

\paragraph{第一组参考题 24}
\begin{xsolution}
由
\[
  f'(x)=\frac1{x^2+f^2(x)}>0
\]
知 $f$ 单调增加，故 $f(x)\geqslant f(1)=1$。于是
\[
  0<f'(x)=\frac1{x^2+f^2(x)}
  \leqslant\frac1{x^2+1}.
\]
因此对 $X>1$，
\[
  1<f(X)
  =1+\int_1^X\frac{\dd x}{x^2+f^2(x)}
  <1+\int_1^{+\infty}\frac{\dd x}{x^2+1}
  =1+\frac\pi4.
\]
所以 $f$ 单调有界，存在有限极限 $L=f(+\infty)$。并且因 $x>1$ 时 $f(x)>1$，有严格不等式
\[
  L-1=\int_1^{+\infty}\frac{\dd x}{x^2+f^2(x)}
  <\int_1^{+\infty}\frac{\dd x}{x^2+1}
  =\frac\pi4.
\]
故
\[
  \boxed{f(+\infty)<1+\frac\pi4}.
\]
\end{xsolution}

\paragraph{第一组参考题 25}
\begin{xsolution}
对 $a\geqslant0$ 定义
\[
  F(a)=\int_0^a\left(\int_x^a\frac{\sin t}{t}\,\dd t\right)\dd x,
\]
其中把 $\sin t/t$ 在 $t=0$ 处连续补定义为 $1$。显然 $F(0)=0$。由变动积分限求导，
\begin{align*}
  F'(a)
  &=\int_0^a\frac{\sin a}{a}\,\dd x\\
  &=\sin a,
\end{align*}
在 $a=0$ 处也按连续性成立。因此
\[
  F(a)=1-\cos a.
\]
取 $a=2\pi$，得到
\[
  \boxed{\int_0^{2\pi}\left(\int_x^{2\pi}\frac{\sin t}{t}\,\dd t\right)\dd x
  =F(2\pi)=0}.
\]
\end{xsolution}


\subsection*{10.5.2 第二组参考题}

\paragraph{第二组参考题 1（1）--（2）}
\begin{xsolution}
先设 $f$ 单调增加。对 $x>0$，由
\[
  f(0+)\leqslant f(t)\leqslant f(x)\qquad(0<t<x)
\]
得到
\[
  f(0+)\leqslant F(x)\leqslant f(x).
\]
若 $0<x<y$，则
\begin{align*}
  yF(y)
  &=xF(x)+\int_x^y f(t)\,\dd t\\
  &\geqslant xF(x)+(y-x)f(x)\\
  &\geqslant xF(x)+(y-x)F(x)=yF(x),
\end{align*}
故 $F(y)\geqslant F(x)$，即 $F$ 也单调增加。若 $f$ 单调减少，上述不等式全部反向，同理得到 $F$ 单调减少。因此 $F$ 与 $f$ 具有相同的单调性。

对任意 $x_0>0$，函数
\[
  x\longmapsto\int_0^x f(t)\,\dd t
\]
连续，所以 $F$ 在 $(0,+\infty)$ 连续。又由单调函数右极限的定义，给定 $\varepsilon>0$，可取 $\delta>0$，使 $0<t<\delta$ 时
\[
  \abs{f(t)-f(0+)}<\varepsilon.
\]
于是 $0<x<\delta$ 时
\[
  \abs{F(x)-f(0+)}
  \leqslant\frac1x\int_0^x\abs{f(t)-f(0+)}\,\dd t<\varepsilon.
\]
因此 $F(x)\to f(0+)=F(0)$，故 $F$ 在 $[0,+\infty)$ 连续。这证明了（1）。

再证（2）。仍先设 $f$ 单调增加，并记扩展实数极限
\[
  L=f(+\infty).
\]
由 $F(x)\leqslant f(x)$，有
\[
  \limsup_{x\to\infty}F(x)\leqslant L.
\]
另一方面，固定 $A>0$。当 $x>A$ 时
\[
  F(x)=\frac1x\int_0^A f(t)\,\dd t
  +\frac1x\int_A^x f(t)\,\dd t
  \geqslant\frac1x\int_0^A f(t)\,\dd t
  +\frac{x-A}{x}f(A).
\]
令 $x\to\infty$，得到
\[
  \liminf_{x\to\infty}F(x)\geqslant f(A).
\]
再令 $A\to\infty$，即有 $\liminf F\geqslant L$。因此
\[
  F(+\infty)=L=f(+\infty).
\]
$f$ 单调减少时对 $-f$ 应用上述结论即可，故（2）在两种单调性下都成立。
\end{xsolution}

\paragraph{第二组参考题 2}
\begin{xsolution}
必要性是直接的。若 $f\in R[a,b]$ 且积分为 $I$，任选满足 $\norm{P_k}\to0$ 的分割序列（例如等距分割），则由 Riemann 积分定义，对任意从属于 $P_k$ 的介点集都有
\[
  \sum_i f(\xi_{k,i})\Delta x_{k,i}\longrightarrow I,
\]
且极限与介点选择无关。

以下证充分性。先注意题设条件本身迫使 $f$ 有界。反设 $f$ 在 $[a,b]$ 上无界，则对每个 $k$，有限个子区间覆盖 $[a,b]$，故 $P_k$ 至少有一个子区间上 $f$ 无界。可在该子区间内选择介点，使这一项的绝对值大到足以令整个第 $k$ 个 Riemann 和的绝对值大于 $k$；其余介点任取。这样便得到一列合法的介点集，其 Riemann 和不可能趋于有限数 $I$，与“极限值不依赖于介点集选取”矛盾。因此 $f$ 在 $[a,b]$ 上有界。

对 $P_k$ 的第 $i$ 个子区间记
\[
  M_{k,i}=\sup f,
  \qquad
  m_{k,i}=\inf f,
  \qquad
  \omega_{k,i}=M_{k,i}-m_{k,i}.
\]
令 $\varepsilon_k\downarrow0$。在每个子区间中分别选择介点 $\alpha_{k,i},\beta_{k,i}$，使
\[
  f(\alpha_{k,i})>M_{k,i}-\varepsilon_k,
  \qquad
  f(\beta_{k,i})<m_{k,i}+\varepsilon_k.
\]
由题设的“极限值不依赖于介点集选取”，有
\[
  \sum_i f(\alpha_{k,i})\Delta x_{k,i}\to I,
  \qquad
  \sum_i f(\beta_{k,i})\Delta x_{k,i}\to I.
\]
两式相减并利用上述逼近，得到
\begin{align*}
  0&\leqslant\sum_i\omega_{k,i}\Delta x_{k,i}\\
  &\leqslant
  \sum_i\bigl(f(\alpha_{k,i})-f(\beta_{k,i})\bigr)\Delta x_{k,i}
  +2\varepsilon_k(b-a)\longrightarrow0.
\end{align*}
因此存在分割使振幅面积任意小。由可积的第二充分必要条件，$f\in R[a,b]$。

既然 $f$ 已经 Riemann 可积，沿给定分割序列的任意介点和必趋于 $\int_a^b f$；题设又说它趋于 $I$，由极限唯一性
\[
  \boxed{\int_a^b f(x)\,\dd x=I}.
\]
\end{xsolution}

\paragraph{第二组参考题 3}
\begin{xsolution}
记右端点和
\[
  R(P)=\sum_{i=1}^n f(x_i)\Delta x_i.
\]
题设说明：当 $\norm P\to0$ 时，对任意分割 $P$，$R(P)$ 一致趋于 $I$。特别地，对任意 $\varepsilon>0$，存在 $\delta>0$，使任意两个细度小于 $\delta$ 的分割 $P,Q$ 都满足
\[
  \abs{R(P)-R(Q)}<2\varepsilon.
\]

下面证明 $f$ 的不连续点集为零测度集。由于题设已经假定 $f$ 有界，可定义一点的振幅
\[
  \omega_f(x)=\lim_{r\to0+}
  \left(\sup_{(x-r,x+r)\cap[a,b]}f-
  \inf_{(x-r,x+r)\cap[a,b]}f\right).
\]
对 $\eta>0$，令
\[
  E_\eta=\{x\in[a,b]\mid\omega_f(x)\geqslant\eta\}.
\]
$E_\eta$ 是闭集：若 $x_n\in E_\eta$ 且 $x_n\to x$，则 $x$ 的任意邻域最终包含 $x_n$ 的某个足够小邻域，故其中 $f$ 的振幅仍至少为 $\eta$。因此 $E_\eta$ 是紧集。

反设某个 $E_\eta$ 不是零测度集。由于
\[
  E_\eta\subseteq\{a,b\}\cup
  \bigcup_{m=1}^\infty
  \left(E_\eta\cap\left[a+\frac1m,b-\frac1m\right]\right),
\]
而有限集与零测度集的可列并仍为零测度集，必存在某个 $m$，使紧集
\[
  K=E_\eta\cap\left[a+\frac1m,b-\frac1m\right]
\]
不是零测度集。因此存在 $\lambda>0$，使 $K$ 的任意开区间覆盖的总长度都不小于 $\lambda$。

先说明一个一维的有限区间选取事实：若有限个开区间覆盖一个点集，则可以从中选出两两不交的子族 $J_1,\ldots,J_N$，使原来所有区间的并包含在 $3J_1\cup\cdots\cup3J_N$ 中。证明是每次选取剩余区间中长度最大的一个，并删去所有与它相交的区间；被删去的区间长度不超过所选区间，又与其相交，因而包含在把所选区间同心扩大三倍所得的区间中。于是若原有限区间族覆盖 $E_\eta\cap(a,b)$，便有
\[
  \sum_{j=1}^N\abs{J_j}\geqslant\frac\lambda3.
\]

现在取一个稍后指定的 $\delta>0$。对每个 $x\in K$，由 $\omega_f(x)\geqslant\eta$，可取任意小的 $r_x>0$，使
\[
  J_x=(x-4r_x,x+4r_x)\subset(a,b),
  \qquad \abs{J_x}<\delta,
\]
并能在中央小区间 $(x-r_x,x+r_x)$ 中找到两点 $y_x<z_x$，满足
\[
  \abs{f(y_x)-f(z_x)}>\frac\eta2.
\]
这些 $J_x$ 构成紧集 $K$ 的开覆盖。由紧致性取有限子覆盖，再用上面的选取事实得到两两不交的 $J_j=(x_j-4r_j,x_j+4r_j)$，且
\[
  \sum_j\abs{J_j}\geqslant\frac\lambda3.
\]
在每个 $J_j$ 中记
\[
  L_j=x_j-2r_j,
\]
并选 $y_j<z_j$ 于 $(x_j-r_j,x_j+r_j)$，使
\[
  \abs{f(y_j)-f(z_j)}>\frac\eta2.
\]
于是
\[
  y_j-L_j>r_j=\frac{\abs{J_j}}8,
  \qquad
  z_j-L_j<3r_j<\delta.
\]

构造一个公共分割 $P_0$，包含所有 $L_j,z_j$，并只在这些互不相交的区间 $[L_j,z_j]$ 之外补充分点，使 $\norm{P_0}<\delta$。从 $P_0$ 构造两个分割 $P_+,P_-$：若 $f(y_j)>f(z_j)$，只把 $y_j$ 插入 $P_+$；若 $f(y_j)<f(z_j)$，只把 $y_j$ 插入 $P_-$。由于各 $[L_j,z_j]$ 两两不交，不同 $j$ 的改动互不影响。若在一个区间 $[L_j,z_j]$ 中插入 $y_j$，右端点和相对于不插入时的变化恰为
\[
  \bigl(f(y_j)-f(z_j)\bigr)(y_j-L_j).
\]
按上述正负号分配后，各项对 $R(P_+)-R(P_-)$ 都作正贡献，因此
\begin{align*}
  R(P_+)-R(P_-)
  &>\sum_j\frac\eta2\cdot\frac{\abs{J_j}}8\\
  &\geqslant\frac{\eta\lambda}{48}.
\end{align*}

另一方面，由题设的统一收敛性，可先取 $\delta$ 足够小，使任意细度小于 $\delta$ 的右端点和都与 $I$ 相差小于 $\eta\lambda/192$。于是
\[
  \abs{R(P_+)-R(P_-)}<\frac{\eta\lambda}{96},
\]
与上面的下界矛盾。故每个 $E_\eta$ 都是零测度集。

函数在一点不连续当且仅当该点振幅大于 $0$，所以
\[
  D(f)=\bigcup_{m=1}^\infty E_{1/m}.
\]
零测度集的可列并仍为零测度集，因此 $D(f)$ 为零测度集。由 Lebesgue 判别定理，$f\in R[a,b]$。

最后，右端点和本来就是 $f$ 的一种 Riemann 和。既然 $f$ 已可积，必有
\[
  R(P)\longrightarrow\int_a^b f(x)\,\dd x
  \qquad(\norm P\to0).
\]
题设又给出同一极限为 $I$，由极限唯一性
\[
  \boxed{\int_a^b f(x)\,\dd x=I}.
\]
\end{xsolution}

\paragraph{第二组参考题 4（Riemann 定理）}
\begin{xsolution}
记周期函数 $g$ 的一个周期平均值为
\[
  \bar g=\frac1T\int_0^Tg(t)\,\dd t.
\]
先证明一个基本估计。由于 $g$ 在 $[0,T]$ 上可积，所以有界，设 $\abs g\leqslant M$。对任意固定有界区间 $[\alpha,\beta]$，令 $u=px$，则
\[
  \int_\alpha^\beta g(px)\,\dd x
  =\frac1p\int_{p\alpha}^{p\beta}g(u)\,\dd u.
\]
把 $[p\alpha,p\beta]$ 分成若干完整的长度为 $T$ 的周期段和至多两个余段。完整周期段每段积分为 $T\bar g$，两个余段总长度小于 $2T$，故存在与 $\alpha,\beta,p$ 无关的常数 $C$，使
\[
  \int_\alpha^\beta g(px)\,\dd x
  =(\beta-\alpha)\bar g+\bigO\left(\frac1p\right),
\]
而且在 $\alpha,\beta\in[a,b]$ 时该 $\bigO(1/p)$ 一致成立。

若 $s$ 是阶梯函数，在有限个子区间上分别取常数 $c_j$，则逐段应用上式即得
\[
  \lim_{p\to+\infty}\int_a^b s(x)g(px)\,\dd x
  =\bar g\int_a^b s(x)\,\dd x.
\]

现在处理一般的 $f\in R[a,b]$。给定 $\varepsilon>0$，由第一组参考题 5（1）可取阶梯函数 $s$，使
\[
  \int_a^b\abs{f(x)-s(x)}\,\dd x<\varepsilon.
\]
于是
\[
  \abs*{\int_a^b(f-s)g(px)\,\dd x}
  \leqslant M\varepsilon,
\]
且
\[
  \abs*{\bar g\int_a^b(f-s)\,\dd x}
  \leqslant\abs{\bar g}\varepsilon.
\]
对固定的 $s$ 令 $p\to\infty$，再令 $\varepsilon\to0$，得到
\[
  \boxed{
  \lim_{p\to+\infty}\int_a^b f(x)g(px)\,\dd x
  =\frac1T\int_0^Tg(t)\,\dd t\int_a^b f(x)\,\dd x}.
\]
\end{xsolution}

\paragraph{第二组参考题 5}
\begin{xsolution}
先由题设正交条件得到一个简单恒等式。写
\[
  f(x)=f(0)+xq(x),
\]
其中 $q$ 是次数不超过 $n-1$ 的多项式。于是 $xq(x)$ 是 $x,x^2,\ldots,x^n$ 的线性组合，故
\[
  \int_0^1 f(x)xq(x)\,\dd x=0.
\]
所以
\begin{equation}
  \int_0^1f^2(x)\,\dd x
  =f(0)\int_0^1f(x)\,\dd x.
  \tag{*}
\end{equation}
只需证明
\[
  f(0)=(n+1)^2\int_0^1f(x)\,\dd x.
\]

构造
\[
  \Phi(x)=x^{n+1}(1-x)^n,
  \qquad
  r_n(x)=\frac1x\frac{\dd^n}{\dd x^n}\Phi(x).
\]
因为 $\Phi^{(n)}(x)$ 含因子 $x$，$r_n$ 是一个 $n$ 次多项式。对 $k=1,\ldots,n$，连续分部积分 $n$ 次；由于 $\Phi$ 在 $0,1$ 两端具有足够高阶零点，所有边界项都为 $0$，而 $(x^{k-1})^{(n)}=0$，故
\[
  \int_0^1x^kr_n(x)\,\dd x
  =\int_0^1x^{k-1}\Phi^{(n)}(x)\,\dd x=0.
\]
这 $n$ 个线性条件彼此独立。事实上，若存在常数 $\lambda_1,\ldots,\lambda_n$，使对每个次数不超过 $n$ 的多项式 $p$ 都有
\[
  \sum_{k=1}^n\lambda_k\int_0^1x^kp(x)\,\dd x=0,
\]
令 $q(x)=\sum_{k=1}^n\lambda_kx^k$ 并取 $p=q$，就得到 $\int_0^1q^2=0$，故 $q\equiv0$，即所有 $\lambda_k=0$。因此这 $n$ 个条件在 $n+1$ 维多项式空间中确实留下一个一维解空间。既然 $r_n$ 是其中的非零元素，必有 $f=C r_n$。

由 $\Phi(x)=x^{n+1}+\bigO(x^{n+2})$（$x\to0$），
\[
  r_n(0)=(n+1)!.
\]
另一方面，连续分部积分 $n$ 次，边界项仍全为 $0$，得到
\begin{align*}
  \int_0^1r_n(x)\,\dd x
  &=\int_0^1\frac{\Phi^{(n)}(x)}x\,\dd x\\
  &=n!\int_0^1(1-x)^n\,\dd x
  =\frac{n!}{n+1}.
\end{align*}
因此
\[
  \frac{f(0)}{\int_0^1f}
  =\frac{(n+1)!}{n!/(n+1)}=(n+1)^2.
\]
代回 $(*)$，即得
\[
  \boxed{\int_0^1f^2(x)\,\dd x
  =(n+1)^2\left(\int_0^1f(x)\,\dd x\right)^2}.
\]
\end{xsolution}

\paragraph{第二组参考题 6（1）}
\begin{xsolution}
用复指数或二项式展开
\[
  \cos^n x=2^{-n}\sum_{k=0}^n C_n^k\ee^{\ii(n-2k)x}.
\]
乘以 $\cos nx$ 后取实部；也可直接使用积化和差公式。只有频率为 $0$ 的常数项在 $[0,\pi/2]$ 上产生非零平均，其他余弦项都是 $\cos(2jx)$，积分均为 $0$。常数项系数为 $2^{-n}$，故
\[
  \boxed{\int_0^{\pi/2}\cos^n x\cos nx\,\dd x
  =\frac{\pi}{2^{n+1}}}.
\]
\end{xsolution}

\paragraph{第二组参考题 6（2）}
\begin{xsolution}
记
\[
  b_n=\int_0^{\pi/2}\cos^n x\sin nx\,\dd x,
  \qquad b_0=0.
\]
对 $n\geqslant1$ 分部积分：
\begin{align*}
  b_n
  &=\frac1n-\int_0^{\pi/2}\cos^{n-1}x\sin x\cos nx\,\dd x.
\end{align*}
利用
\[
  \sin x\cos nx=\sin nx\cos x-\sin(n-1)x,
\]
可得
\[
  \int_0^{\pi/2}\cos^{n-1}x\sin x\cos nx\,\dd x
  =b_n-b_{n-1}.
\]
故
\[
  2b_n=\frac1n+b_{n-1},
  \qquad
  b_n=\frac1{2n}+\frac12b_{n-1}.
\]
递推即得
\[
  \boxed{
  \int_0^{\pi/2}\cos^n x\sin nx\,\dd x
  =\frac1{2^{n+1}}\sum_{k=1}^n\frac{2^k}{k}}.
\]
\end{xsolution}

\paragraph{第二组参考题 7}
\begin{xsolution}
只需估计绝对值：
\[
  \abs*{\int_0^1\cos^n\frac1x\,\dd x}
  \leqslant\int_0^1\abs*{\cos\frac1x}^{\,n}\,\dd x.
\]
给定 $\varepsilon>0$，先取 $\delta>0$ 使 $\delta<\varepsilon/3$。则
\[
  \int_0^\delta\abs*{\cos\frac1x}^{\,n}\,\dd x\leqslant\delta<\frac\varepsilon3.
\]
在紧区间 $[\delta,1]$ 上，集合
\[
  E=\left\{x:\abs*{\cos\frac1x}=1\right\}
\]
只有有限个点，因为它等价于 $1/x=k\pi$。用有限个小开区间覆盖这些点，并使覆盖总长度小于 $\varepsilon/3$。在其余紧集 $K$ 上，连续函数 $\abs{\cos(1/x)}$ 的最大值满足
\[
  q:=\max_K\abs*{\cos\frac1x}<1.
\]
于是
\[
  \int_\delta^1\abs*{\cos\frac1x}^{\,n}\,\dd x
  \leqslant\frac\varepsilon3+(1-\delta)q^n.
\]
当 $n$ 充分大时最后一项小于 $\varepsilon/3$。因此整个绝对值小于 $\varepsilon$，从而
\[
  \boxed{\lim_{n\to\infty}\int_0^1\cos^n\frac1x\,\dd x=0}.
\]
\end{xsolution}

\paragraph{第二组参考题 8}
\begin{xsolution}
先作代数分解：
\[
  \frac{4\sqrt2-8x^3-4\sqrt2x^4-8x^5}{1-x^8}
  =-\frac{4(x-\sqrt2)}{x^2-\sqrt2x+1}
  +\frac2{x+1}+\frac2{x-1}.
\]
在 $0\leqslant x\leqslant1/\sqrt2<1$ 上各项均无奇点。注意
\[
  x^2-\sqrt2x+1
  =\frac12\bigl((\sqrt2x-1)^2+1\bigr).
\]
直接积分可得一个原函数
\[
  A(x)=2\ln(1-x^2)
  -2\ln(x^2-\sqrt2x+1)
  +4\arctan(\sqrt2x-1).
\]
在 $x=1/\sqrt2$ 时，前两个对数中的真数都等于 $1/2$，且反正切项为 $0$，故
\[
  A(1/\sqrt2)=0.
\]
在 $x=0$ 时两个对数项为 $0$，而
\[
  A(0)=4\arctan(-1)=-\pi.
\]
所以
\[
  \boxed{
  \int_0^{1/\sqrt2}
  \frac{4\sqrt2-8x^3-4\sqrt2x^4-8x^5}{1-x^8}\,\dd x
  =A(1/\sqrt2)-A(0)=\pi}.
\]
\end{xsolution}
